
\documentclass{article}
\usepackage{colortbl}
\usepackage{makecell}
\usepackage{multirow}
\usepackage{supertabular}

\begin{document}

\newcounter{utterance}

\centering \large Interaction Transcript for game `clean\_up', experiment `1\_medium\_7obj\_cs', episode 3 with deepseek{-}v4{-}pro.
\vspace{24pt}

{ \footnotesize  \setcounter{utterance}{1}
\setlength{\tabcolsep}{0pt}
\begin{supertabular}{c@{$\;$}|p{.15\linewidth}@{}p{.15\linewidth}p{.15\linewidth}p{.15\linewidth}p{.15\linewidth}p{.15\linewidth}}
    \# & $\;$A & \multicolumn{4}{c}{Game Master} & $\;\:$B\\
    \hline

    \theutterance \stepcounter{utterance}  
    & & \multicolumn{4}{p{0.6\linewidth}}{
        \cellcolor[rgb]{0.9,0.9,0.9}{
            \makecell[{{p{\linewidth}}}]{
                \texttt{\tiny{[P1$\langle$GM]}}
                \texttt{Jsem váš herní mistr a hrajete společnou hru s následující mřížkou jako herním plánem:} \\
\\ 
\texttt{```} \\
\texttt{ 1234567} \\
\texttt{╔═╤═════╗} \\
\texttt{║◌│◌A◌◌◌║ 1} \\
\texttt{╟─┤◌T◌◌◌║ 2} \\
\texttt{║O└─┐◌◌◌║ 3} \\
\texttt{║◌◌◌├─┐◌║ 4} \\
\texttt{║◌◌S│◌│◌║ 5} \\
\texttt{╟─┬─┴─┼─╢ 6} \\
\texttt{║P│◌DM│◌║ 7} \\
\texttt{╚═╧═══╧═╝} \\
\\ 
\texttt{```} \\
\\ 
\texttt{* Horní okraj zobrazuje souřadnice x rostoucí doprava a pravý okraj souřadnice y rostoucí dolů.} \\
\texttt{* Na vaší mřížce jsou náhodně umístěny následující objekty: 'P', 'O', 'T', 'S', 'D', 'A', 'M'.} \\
\\ 
\\ 
\texttt{Druhý hráč vidí variantu herní desky, kde jsou objekty umístěny na různých náhodných místech. Vy nevidíte desku druhého hráče a on nevidí tu vaši.} \\
\\ 
\\ 
\texttt{**Cíl:**} \\
\\ 
\texttt{Oba hráči musí přesunout objekty na svém pozadí tak, aby stejné objekty skončily na stejných souřadnicích. Musíte komunikovat s druhým hráčem, abyste se dohodli na společné cílové konfiguraci.} \\
\\ 
\\ 
\texttt{**Pravidla:**} \\
\\ 
\texttt{* V každém tahu můžete poslat přesně jeden z následujících dvou příkazů:} \\
\texttt{ 1. `ŘEKNI: <ZPRÁVA>`: odeslání zprávy (vše až do dalšího zalomení řádku) druhému hráči. Přepošlu ji vašemu partnerovi.} \\
\texttt{2. `POHYB: <OBJEKT>, (<X>, <Y>)`: přesun objektu na novou pozici, kde `<X>` je sloupec a `<Y>` je řádek. Budu vás informovat, zda byl váš přesun úspěšný, nebo ne.} \\
\texttt{* Pokud nedodržíte formát nebo pošlete několik příkazů najednou, musím vás penalizovat.} \\
\texttt{* Pokud oba hráči nasbírají více než 16 penalizací, oba prohrajete hru.} \\
\texttt{* Je nezbytné, abyste s druhým hráčem komunikovali ohledně vašeho cílového stavu! *Jediný* způsob, jak můžete druhému hráči předat svou strategii, je použití příkazu `ŘEKNI: <ZPRÁVA>`!} \\
\\ 
\\ 
\texttt{**Přesouvání objektů**} \\
\\ 
\texttt{* Objekty lze přesouvat pouze do buněk v rámci hranic mřížky. Cílová buňka musí být prázdná, tj. musí obsahovat pouze symbol '◌'.} \\
\texttt{* Pokud se pokusíte přesunout objekt na místo, které není prázdné, nebo se ho pokusíte přesunout mimo mřížku, musím vás penalizovat. Dostanete další pokus.} \\
\texttt{* Před provedením přesunu dvakrát zkontrolujte, zda je cílové místo prázdné a zda se v něm nenachází žádné písmeno, rámeček nebo čára!} \\
\\ 
\\ 
\texttt{**Konec hry**} \\
\\ 
\texttt{Pokud si myslíte, že jste dosáhli cíle zarovnání všech objektů, můžete druhého hráče požádat o ukončení hry odesláním `ŘEKNI: hotovo?`. Pokud vás druhý hráč požádá o ukončení hry a vy odpovíte `ŘEKNI: hotovo!`, hra skončí.} \\
\\ 
\texttt{Oba hráči vyhrají, pokud hra skončí během 28 kol, přičemž jedno kolo je definováno jako odeslání platného příkazu oběma hráči.} \\
\\ 
\texttt{**Bodování:**} \\
\\ 
\texttt{Čím blíže jsou stejné objekty na obou herních deskách, tím více bodů získáte. Penalizace vám body snižují. Podaří se vám překonat rekord?} \\
\\ 
\\ 
\texttt{Pošlete prosím zprávu druhému hráči a začněte hru!} \\
            }
        }
    }
    & \\ \\

    \theutterance \stepcounter{utterance}  
    & \multicolumn{4}{p{0.6\linewidth}}{
        \cellcolor[rgb]{0.9,0.9,0.9}{
            \makecell[{{p{\linewidth}}}]{
                \texttt{\tiny{[P1$\rangle$GM]}}
                \texttt{ŘEKNI: Ahoj, jaké máš rozložení objektů na své desce? Tady je moje: P (2,7), O (1,3), T (3,2), S (3,5), D (4,7), A (3,1), M (5,7). Jaká je tvoje konfigurace?} \\
            }
        }
    }
    & & \\ \\

    \theutterance \stepcounter{utterance}  
    & & & \multicolumn{2}{p{0.3\linewidth}}{
        \cellcolor[rgb]{0.9,0.9,0.9}{
            \makecell[{{p{\linewidth}}}]{
                \texttt{\tiny{[GM$|$GM]}}
                \texttt{ŘEKNI: Ahoj, jaké máš rozložení objektů na své desce? Tady je moje: P (2,7), O (1,3), T (3,2), S (3,5), D (4,7), A (3,1), M (5,7). Jaká je tvoje konfigurace?} \\
            }
        }
    }
    & & \\ \\

    \theutterance \stepcounter{utterance}  
    & & \multicolumn{4}{p{0.6\linewidth}}{
        \cellcolor[rgb]{0.9,0.9,0.9}{
            \makecell[{{p{\linewidth}}}]{
                \texttt{\tiny{[GM$\rangle$P2]}}
                \texttt{Jsem váš herní mistr a hrajete společnou hru s následující mřížkou jako herním plánem:} \\
\\ 
\texttt{```} \\
\texttt{ 1234567} \\
\texttt{╔═╤═════╗} \\
\texttt{║P│S◌D◌◌║ 1} \\
\texttt{╟─┤◌◌◌◌◌║ 2} \\
\texttt{║◌└─┐◌A◌║ 3} \\
\texttt{║◌◌M├─┐◌║ 4} \\
\texttt{║◌T◌│◌│O║ 5} \\
\texttt{╟─┬─┴─┼─╢ 6} \\
\texttt{║◌│◌◌◌│◌║ 7} \\
\texttt{╚═╧═══╧═╝} \\
\\ 
\texttt{```} \\
\\ 
\texttt{* Horní okraj zobrazuje souřadnice x rostoucí doprava a pravý okraj souřadnice y rostoucí dolů.} \\
\texttt{* Na vaší mřížce jsou náhodně umístěny následující objekty: 'P', 'O', 'T', 'S', 'D', 'A', 'M'.} \\
\\ 
\\ 
\texttt{Druhý hráč vidí variantu herní desky, kde jsou objekty umístěny na různých náhodných místech. Vy nevidíte desku druhého hráče a on nevidí tu vaši.} \\
\\ 
\\ 
\texttt{**Cíl:**} \\
\\ 
\texttt{Oba hráči musí přesunout objekty na svém pozadí tak, aby stejné objekty skončily na stejných souřadnicích. Musíte komunikovat s druhým hráčem, abyste se dohodli na společné cílové konfiguraci.} \\
\\ 
\\ 
\texttt{**Pravidla:**} \\
\\ 
\texttt{* V každém tahu můžete poslat přesně jeden z následujících dvou příkazů:} \\
\texttt{ 1. `ŘEKNI: <ZPRÁVA>`: odeslání zprávy (vše až do dalšího zalomení řádku) druhému hráči. Přepošlu ji vašemu partnerovi.} \\
\texttt{2. `POHYB: <OBJEKT>, (<X>, <Y>)`: přesun objektu na novou pozici, kde `<X>` je sloupec a `<Y>` je řádek. Budu vás informovat, zda byl váš přesun úspěšný, nebo ne.} \\
\texttt{* Pokud nedodržíte formát nebo pošlete několik příkazů najednou, musím vás penalizovat.} \\
\texttt{* Pokud oba hráči nasbírají více než 16 penalizací, oba prohrajete hru.} \\
\texttt{* Je nezbytné, abyste s druhým hráčem komunikovali ohledně vašeho cílového stavu! *Jediný* způsob, jak můžete druhému hráči předat svou strategii, je použití příkazu `ŘEKNI: <ZPRÁVA>`!} \\
\\ 
\\ 
\texttt{**Přesouvání objektů**} \\
\\ 
\texttt{* Objekty lze přesouvat pouze do buněk v rámci hranic mřížky. Cílová buňka musí být prázdná, tj. musí obsahovat pouze symbol '◌'.} \\
\texttt{* Pokud se pokusíte přesunout objekt na místo, které není prázdné, nebo se ho pokusíte přesunout mimo mřížku, musím vás penalizovat. Dostanete další pokus.} \\
\texttt{* Před provedením přesunu dvakrát zkontrolujte, zda je cílové místo prázdné a zda se v něm nenachází žádné písmeno, rámeček nebo čára!} \\
\\ 
\\ 
\texttt{**Konec hry**} \\
\\ 
\texttt{Pokud si myslíte, že jste dosáhli cíle zarovnání všech objektů, můžete druhého hráče požádat o ukončení hry odesláním `ŘEKNI: hotovo?`. Pokud vás druhý hráč požádá o ukončení hry a vy odpovíte `ŘEKNI: hotovo!`, hra skončí.} \\
\\ 
\texttt{Oba hráči vyhrají, pokud hra skončí během 28 kol, přičemž jedno kolo je definováno jako odeslání platného příkazu oběma hráči.} \\
\\ 
\texttt{**Bodování:**} \\
\\ 
\texttt{Čím blíže jsou stejné objekty na obou herních deskách, tím více bodů získáte. Penalizace vám body snižují. Podaří se vám překonat rekord?} \\
\\ 
\texttt{Druhý hráč zahájil hru odesláním této zprávy:} \\
\\ 
\texttt{"Ahoj, jaké máš rozložení objektů na své desce? Tady je moje: P (2,7), O (1,3), T (3,2), S (3,5), D (4,7), A (3,1), M (5,7). Jaká je tvoje konfigurace?"} \\
\\ 
\texttt{Jaký je váš první příkaz?} \\
            }
        }
    }
    & \\ \\

    \theutterance \stepcounter{utterance}  
    & & & \multicolumn{4}{p{0.6\linewidth}}{
        \cellcolor[rgb]{0.9,0.9,0.9}{
            \makecell[{{p{\linewidth}}}]{
                \texttt{\tiny{[GM$\langle$P2]}}
                \texttt{Ahoj, díky za tvou konfiguraci. Tady je moje:} \\
\\ 
\texttt{P (1,1), S (2,1), D (6,1), A (6,3), M (3,4), T (3,5), O (7,5)} \\
\\ 
\texttt{Teď se musíme dohodnout na společných cílových pozicích. Navrhuji, abychom použili pozice, které jsou co nejblíže průměru našich aktuálních pozic, nebo se dohodli na konkrétním rozložení. Co myslíš?} \\
\\ 
\texttt{Můj první příkaz:} \\
\\ 
\texttt{ŘEKNI: Ahoj, tady je moje konfigurace: P (1,1), S (2,1), D (6,1), A (6,3), M (3,4), T (3,5), O (7,5). Jak chceme postupovat při výběru cílových pozic?} \\
            }
        }
    }
     \\ \\

    \theutterance \stepcounter{utterance}  
    & & & \multicolumn{2}{p{0.3\linewidth}}{
        \cellcolor[rgb]{0.9,0.9,0.9}{
            \makecell[{{p{\linewidth}}}]{
                \texttt{\tiny{[GM$|$GM]}}
                \texttt{Ahoj, díky za tvou konfiguraci. Tady je moje:} \\
\\ 
\texttt{P (1,1), S (2,1), D (6,1), A (6,3), M (3,4), T (3,5), O (7,5)} \\
\\ 
\texttt{Teď se musíme dohodnout na společných cílových pozicích. Navrhuji, abychom použili pozice, které jsou co nejblíže průměru našich aktuálních pozic, nebo se dohodli na konkrétním rozložení. Co myslíš?} \\
\\ 
\texttt{Můj první příkaz:} \\
\\ 
\texttt{ŘEKNI: Ahoj, tady je moje konfigurace: P (1,1), S (2,1), D (6,1), A (6,3), M (3,4), T (3,5), O (7,5). Jak chceme postupovat při výběru cílových pozic?} \\
            }
        }
    }
    & & \\ \\

    \theutterance \stepcounter{utterance}  
    & & & \multicolumn{2}{p{0.3\linewidth}}{
        \cellcolor[rgb]{0.9,0.9,0.9}{
            \makecell[{{p{\linewidth}}}]{
                \texttt{\tiny{[GM$|$GM]}}
                \texttt{Invalid format: head is not empty: 'Ahoj, díky za tvou konfiguraci. Tady je moje:} \\
\\ 
\texttt{P (1,1), S (2,1), D (6,1), A (6,3), M (3,4), T (3,5), O (7,5)} \\
\\ 
\texttt{Teď se musíme dohodnout na společných cílových pozicích. Navrhuji, abychom použili pozice, které jsou co nejblíže průměru našich aktuálních pozic, nebo se dohodli na konkrétním rozložení. Co myslíš?} \\
\\ 
\texttt{Můj první příkaz:} \\
\\ 
\texttt{'} \\
            }
        }
    }
    & & \\ \\

    \theutterance \stepcounter{utterance}  
    & & \multicolumn{4}{p{0.6\linewidth}}{
        \cellcolor[rgb]{0.9,0.9,0.9}{
            \makecell[{{p{\linewidth}}}]{
                \texttt{\tiny{[GM$\rangle$P2]}}
                \texttt{Trest: Vaše zpráva nesmí obsahovat nic před příkazem!} \\
\\ 
\texttt{Ujistěte se, že vaše odpověď obsahuje pouze `ŘEKNI: <ZPRÁVA>` nebo `POHYB: <OBJEKT>, (<X>, <Y>)` a nic jiného!} \\
\\ 
\texttt{Dohromady jste nasbírali 1 z 16 pokut.} \\
\texttt{Zkuste to prosím znovu!} \\
            }
        }
    }
    & \\ \\

    \theutterance \stepcounter{utterance}  
    & & & \multicolumn{4}{p{0.6\linewidth}}{
        \cellcolor[rgb]{0.9,0.9,0.9}{
            \makecell[{{p{\linewidth}}}]{
                \texttt{\tiny{[GM$\langle$P2]}}
                \texttt{ŘEKNI: Ahoj, tady je moje konfigurace: P (1,1), S (2,1), D (6,1), A (6,3), M (3,4), T (3,5), O (7,5). Jak chceme postupovat při výběru cílových pozic?} \\
            }
        }
    }
     \\ \\

    \theutterance \stepcounter{utterance}  
    & & & \multicolumn{2}{p{0.3\linewidth}}{
        \cellcolor[rgb]{0.9,0.9,0.9}{
            \makecell[{{p{\linewidth}}}]{
                \texttt{\tiny{[GM$|$GM]}}
                \texttt{ŘEKNI: Ahoj, tady je moje konfigurace: P (1,1), S (2,1), D (6,1), A (6,3), M (3,4), T (3,5), O (7,5). Jak chceme postupovat při výběru cílových pozic?} \\
            }
        }
    }
    & & \\ \\

    \theutterance \stepcounter{utterance}  
    & & \multicolumn{4}{p{0.6\linewidth}}{
        \cellcolor[rgb]{0.9,0.9,0.9}{
            \makecell[{{p{\linewidth}}}]{
                \texttt{\tiny{[P1$\langle$GM]}}
                \texttt{Vaše zpráva byla předána druhému hráči.} \\
\\ 
\\ 
\texttt{Právě hrajete kolo 1 z maximálně 28.} \\
\texttt{Dohromady jste nasbírali 1 z 16 pokut.} \\
\texttt{Druhý hráč poslal tuto zprávu:} \\
\\ 
\\ 
\texttt{"Ahoj, tady je moje konfigurace: P (1,1), S (2,1), D (6,1), A (6,3), M (3,4), T (3,5), O (7,5). Jak chceme postupovat při výběru cílových pozic?"} \\
\\ 
\\ 
\texttt{Jaký je váš další příkaz?} \\
            }
        }
    }
    & \\ \\

    \theutterance \stepcounter{utterance}  
    & \multicolumn{4}{p{0.6\linewidth}}{
        \cellcolor[rgb]{0.9,0.9,0.9}{
            \makecell[{{p{\linewidth}}}]{
                \texttt{\tiny{[P1$\rangle$GM]}}
                \texttt{ŘEKNI: Super, díky za info! Porovnal jsem naše pozice – každý objekt je na úplně jiném místě, takže se musíme dohodnout na společných cílových souřadnicích.} \\
\\ 
\texttt{Navrhuji, abychom pro každý objekt vybrali cílovou pozici, která je pro oba snadno dosažitelná a vyhýbá se kolizím. Tady je můj návrh cílových pozic:} \\
\\ 
\texttt{P: (1,1) – ty už ho tam máš, já ho přesunu} \\
\texttt{O: (1,3) – já už ho tam mám} \\
\texttt{T: (3,3) – volné u obou?} \\
\texttt{S: (5,5) – nebo jiný návrh?} \\
\texttt{D: (6,6) – ?} \\
\texttt{A: (7,1) – ?} \\
\texttt{M: (5,7) – já ho tam mám} \\
\\ 
\texttt{Co ty na to? Které pozice jsou u tebe volné?} \\
            }
        }
    }
    & & \\ \\

    \theutterance \stepcounter{utterance}  
    & & & \multicolumn{2}{p{0.3\linewidth}}{
        \cellcolor[rgb]{0.9,0.9,0.9}{
            \makecell[{{p{\linewidth}}}]{
                \texttt{\tiny{[GM$|$GM]}}
                \texttt{ŘEKNI: Super, díky za info! Porovnal jsem naše pozice – každý objekt je na úplně jiném místě, takže se musíme dohodnout na společných cílových souřadnicích.} \\
\\ 
\texttt{Navrhuji, abychom pro každý objekt vybrali cílovou pozici, která je pro oba snadno dosažitelná a vyhýbá se kolizím. Tady je můj návrh cílových pozic:} \\
\\ 
\texttt{P: (1,1) – ty už ho tam máš, já ho přesunu} \\
\texttt{O: (1,3) – já už ho tam mám} \\
\texttt{T: (3,3) – volné u obou?} \\
\texttt{S: (5,5) – nebo jiný návrh?} \\
\texttt{D: (6,6) – ?} \\
\texttt{A: (7,1) – ?} \\
\texttt{M: (5,7) – já ho tam mám} \\
\\ 
\texttt{Co ty na to? Které pozice jsou u tebe volné?} \\
            }
        }
    }
    & & \\ \\

    \theutterance \stepcounter{utterance}  
    & & & \multicolumn{2}{p{0.3\linewidth}}{
        \cellcolor[rgb]{0.9,0.9,0.9}{
            \makecell[{{p{\linewidth}}}]{
                \texttt{\tiny{[GM$|$GM]}}
                \texttt{Invalid format: tail is not empty: '} \\
\\ 
\texttt{Navrhuji, abychom pro každý objekt vybrali cílovou pozici, která je pro oba snadno dosažitelná a vyhýbá se kolizím. Tady je můj návrh cílových pozic:} \\
\\ 
\texttt{P: (1,1) – ty už ho tam máš, já ho přesunu} \\
\texttt{O: (1,3) – já už ho tam mám} \\
\texttt{T: (3,3) – volné u obou?} \\
\texttt{S: (5,5) – nebo jiný návrh?} \\
\texttt{D: (6,6) – ?} \\
\texttt{A: (7,1) – ?} \\
\texttt{M: (5,7) – já ho tam mám} \\
\\ 
\texttt{Co ty na to? Které pozice jsou u tebe volné?} \\
            }
        }
    }
    & & \\ \\

    \theutterance \stepcounter{utterance}  
    & & \multicolumn{4}{p{0.6\linewidth}}{
        \cellcolor[rgb]{0.9,0.9,0.9}{
            \makecell[{{p{\linewidth}}}]{
                \texttt{\tiny{[P1$\langle$GM]}}
                \texttt{Trest: Vaše zpráva nesmí obsahovat nic po příkazu!} \\
\\ 
\texttt{Ujistěte se, že vaše odpověď obsahuje pouze `ŘEKNI: <ZPRÁVA>` nebo `POHYB: <OBJEKT>, (<X>, <Y>)` a nic jiného!} \\
\\ 
\texttt{Dohromady jste nasbírali 2 z 16 pokut.} \\
\texttt{Zkuste to prosím znovu!} \\
            }
        }
    }
    & \\ \\

    \theutterance \stepcounter{utterance}  
    & \multicolumn{4}{p{0.6\linewidth}}{
        \cellcolor[rgb]{0.9,0.9,0.9}{
            \makecell[{{p{\linewidth}}}]{
                \texttt{\tiny{[P1$\rangle$GM]}}
                \texttt{ŘEKNI: Super, díky za info! Porovnal jsem naše pozice a navrhuji cílové souřadnice: P (1,1), O (1,3), T (3,3), S (5,5), D (6,6), A (7,1), M (5,7). Které pozice jsou u tebe volné?} \\
            }
        }
    }
    & & \\ \\

    \theutterance \stepcounter{utterance}  
    & & & \multicolumn{2}{p{0.3\linewidth}}{
        \cellcolor[rgb]{0.9,0.9,0.9}{
            \makecell[{{p{\linewidth}}}]{
                \texttt{\tiny{[GM$|$GM]}}
                \texttt{ŘEKNI: Super, díky za info! Porovnal jsem naše pozice a navrhuji cílové souřadnice: P (1,1), O (1,3), T (3,3), S (5,5), D (6,6), A (7,1), M (5,7). Které pozice jsou u tebe volné?} \\
            }
        }
    }
    & & \\ \\

    \theutterance \stepcounter{utterance}  
    & & \multicolumn{4}{p{0.6\linewidth}}{
        \cellcolor[rgb]{0.9,0.9,0.9}{
            \makecell[{{p{\linewidth}}}]{
                \texttt{\tiny{[GM$\rangle$P2]}}
                \texttt{Vaše zpráva byla předána druhému hráči.} \\
\\ 
\\ 
\texttt{Právě hrajete kolo 2 z maximálně 28.} \\
\texttt{Dohromady jste nasbírali 2 z 16 pokut.} \\
\texttt{Druhý hráč poslal tuto zprávu:} \\
\\ 
\\ 
\texttt{"Super, díky za info! Porovnal jsem naše pozice a navrhuji cílové souřadnice: P (1,1), O (1,3), T (3,3), S (5,5), D (6,6), A (7,1), M (5,7). Které pozice jsou u tebe volné?"} \\
\\ 
\\ 
\texttt{Jaký je váš další příkaz?} \\
            }
        }
    }
    & \\ \\

    \theutterance \stepcounter{utterance}  
    & & & \multicolumn{4}{p{0.6\linewidth}}{
        \cellcolor[rgb]{0.9,0.9,0.9}{
            \makecell[{{p{\linewidth}}}]{
                \texttt{\tiny{[GM$\langle$P2]}}
                \texttt{ŘEKNI: Zkontroloval jsem navržené cílové pozice. U mě jsou volné: (1,3), (3,3), (5,5), (6,6), (7,1), (5,7). P (1,1) už sedí, S (2,1) a D (6,1) a A (6,3) a M (3,4) a T (3,5) a O (7,5) jsou obsazené, takže je musím přesunout. Souhlasím s tvým návrhem, začnu přesouvat. Můžeš potvrdit, že souhlasíš s cílovými pozicemi?} \\
            }
        }
    }
     \\ \\

    \theutterance \stepcounter{utterance}  
    & & & \multicolumn{2}{p{0.3\linewidth}}{
        \cellcolor[rgb]{0.9,0.9,0.9}{
            \makecell[{{p{\linewidth}}}]{
                \texttt{\tiny{[GM$|$GM]}}
                \texttt{ŘEKNI: Zkontroloval jsem navržené cílové pozice. U mě jsou volné: (1,3), (3,3), (5,5), (6,6), (7,1), (5,7). P (1,1) už sedí, S (2,1) a D (6,1) a A (6,3) a M (3,4) a T (3,5) a O (7,5) jsou obsazené, takže je musím přesunout. Souhlasím s tvým návrhem, začnu přesouvat. Můžeš potvrdit, že souhlasíš s cílovými pozicemi?} \\
            }
        }
    }
    & & \\ \\

    \theutterance \stepcounter{utterance}  
    & & \multicolumn{4}{p{0.6\linewidth}}{
        \cellcolor[rgb]{0.9,0.9,0.9}{
            \makecell[{{p{\linewidth}}}]{
                \texttt{\tiny{[P1$\langle$GM]}}
                \texttt{Vaše zpráva byla předána druhému hráči.} \\
\\ 
\\ 
\texttt{Právě hrajete kolo 2 z maximálně 28.} \\
\texttt{Dohromady jste nasbírali 2 z 16 pokut.} \\
\texttt{Druhý hráč poslal tuto zprávu:} \\
\\ 
\\ 
\texttt{"Zkontroloval jsem navržené cílové pozice. U mě jsou volné: (1,3), (3,3), (5,5), (6,6), (7,1), (5,7). P (1,1) už sedí, S (2,1) a D (6,1) a A (6,3) a M (3,4) a T (3,5) a O (7,5) jsou obsazené, takže je musím přesunout. Souhlasím s tvým návrhem, začnu přesouvat. Můžeš potvrdit, že souhlasíš s cílovými pozicemi?"} \\
\\ 
\\ 
\texttt{Jaký je váš další příkaz?} \\
            }
        }
    }
    & \\ \\

    \theutterance \stepcounter{utterance}  
    & \multicolumn{4}{p{0.6\linewidth}}{
        \cellcolor[rgb]{0.9,0.9,0.9}{
            \makecell[{{p{\linewidth}}}]{
                \texttt{\tiny{[P1$\rangle$GM]}}
                \texttt{ŘEKNI: Potvrzuji souhlas s cílovými pozicemi: P (1,1), O (1,3), T (3,3), S (5,5), D (6,6), A (7,1), M (5,7). Začínám s přesuny.} \\
            }
        }
    }
    & & \\ \\

    \theutterance \stepcounter{utterance}  
    & & & \multicolumn{2}{p{0.3\linewidth}}{
        \cellcolor[rgb]{0.9,0.9,0.9}{
            \makecell[{{p{\linewidth}}}]{
                \texttt{\tiny{[GM$|$GM]}}
                \texttt{ŘEKNI: Potvrzuji souhlas s cílovými pozicemi: P (1,1), O (1,3), T (3,3), S (5,5), D (6,6), A (7,1), M (5,7). Začínám s přesuny.} \\
            }
        }
    }
    & & \\ \\

    \theutterance \stepcounter{utterance}  
    & & \multicolumn{4}{p{0.6\linewidth}}{
        \cellcolor[rgb]{0.9,0.9,0.9}{
            \makecell[{{p{\linewidth}}}]{
                \texttt{\tiny{[GM$\rangle$P2]}}
                \texttt{Vaše zpráva byla předána druhému hráči.} \\
\\ 
\\ 
\texttt{Právě hrajete kolo 3 z maximálně 28.} \\
\texttt{Dohromady jste nasbírali 2 z 16 pokut.} \\
\texttt{Druhý hráč poslal tuto zprávu:} \\
\\ 
\\ 
\texttt{"Potvrzuji souhlas s cílovými pozicemi: P (1,1), O (1,3), T (3,3), S (5,5), D (6,6), A (7,1), M (5,7). Začínám s přesuny."} \\
\\ 
\\ 
\texttt{Jaký je váš další příkaz?} \\
            }
        }
    }
    & \\ \\

    \theutterance \stepcounter{utterance}  
    & & & \multicolumn{4}{p{0.6\linewidth}}{
        \cellcolor[rgb]{0.9,0.9,0.9}{
            \makecell[{{p{\linewidth}}}]{
                \texttt{\tiny{[GM$\langle$P2]}}
                \texttt{POHYB: S, (5,5)} \\
            }
        }
    }
     \\ \\

    \theutterance \stepcounter{utterance}  
    & & & \multicolumn{2}{p{0.3\linewidth}}{
        \cellcolor[rgb]{0.9,0.9,0.9}{
            \makecell[{{p{\linewidth}}}]{
                \texttt{\tiny{[GM$|$GM]}}
                \texttt{POHYB: S, (5,5)} \\
            }
        }
    }
    & & \\ \\

    \theutterance \stepcounter{utterance}  
    & & & \multicolumn{2}{p{0.3\linewidth}}{
        \cellcolor[rgb]{0.9,0.9,0.9}{
            \makecell[{{p{\linewidth}}}]{
                \texttt{\tiny{[GM$|$GM]}}
                \texttt{S moved to (5, 5)} \\
            }
        }
    }
    & & \\ \\

    \theutterance \stepcounter{utterance}  
    & & \multicolumn{4}{p{0.6\linewidth}}{
        \cellcolor[rgb]{0.9,0.9,0.9}{
            \makecell[{{p{\linewidth}}}]{
                \texttt{\tiny{[P1$\langle$GM]}}
                \texttt{Vaše zpráva byla předána druhému hráči.} \\
\\ 
\\ 
\texttt{Právě hrajete kolo 3 z maximálně 28.} \\
\texttt{Dohromady jste nasbírali 2 z 16 pokut.} \\
\texttt{Druhý hráč přesunul objekt na své mřížce.} \\
\\ 
\\ 
\texttt{Jaký je váš další příkaz?} \\
            }
        }
    }
    & \\ \\

    \theutterance \stepcounter{utterance}  
    & \multicolumn{4}{p{0.6\linewidth}}{
        \cellcolor[rgb]{0.9,0.9,0.9}{
            \makecell[{{p{\linewidth}}}]{
                \texttt{\tiny{[P1$\rangle$GM]}}
                \texttt{POHYB: P, (1,1)} \\
            }
        }
    }
    & & \\ \\

    \theutterance \stepcounter{utterance}  
    & & & \multicolumn{2}{p{0.3\linewidth}}{
        \cellcolor[rgb]{0.9,0.9,0.9}{
            \makecell[{{p{\linewidth}}}]{
                \texttt{\tiny{[GM$|$GM]}}
                \texttt{POHYB: P, (1,1)} \\
            }
        }
    }
    & & \\ \\

    \theutterance \stepcounter{utterance}  
    & & & \multicolumn{2}{p{0.3\linewidth}}{
        \cellcolor[rgb]{0.9,0.9,0.9}{
            \makecell[{{p{\linewidth}}}]{
                \texttt{\tiny{[GM$|$GM]}}
                \texttt{P moved to (1, 1)} \\
            }
        }
    }
    & & \\ \\

    \theutterance \stepcounter{utterance}  
    & & \multicolumn{4}{p{0.6\linewidth}}{
        \cellcolor[rgb]{0.9,0.9,0.9}{
            \makecell[{{p{\linewidth}}}]{
                \texttt{\tiny{[GM$\rangle$P2]}}
                \texttt{Úspěšně přesunuto 'S' do (5,5). Vaše aktualizovaná mřížka vypadá takto:} \\
\texttt{```} \\
\texttt{ 1234567} \\
\texttt{╔═╤═════╗} \\
\texttt{║P│◌◌D◌◌║ 1} \\
\texttt{╟─┤◌◌◌◌◌║ 2} \\
\texttt{║◌└─┐◌A◌║ 3} \\
\texttt{║◌◌M├─┐◌║ 4} \\
\texttt{║◌T◌│S│O║ 5} \\
\texttt{╟─┬─┴─┼─╢ 6} \\
\texttt{║◌│◌◌◌│◌║ 7} \\
\texttt{╚═╧═══╧═╝} \\
\texttt{```} \\
\texttt{Právě hrajete kolo 4 z maximálně 28.} \\
\texttt{Dohromady jste nasbírali 2 z 16 pokut.} \\
\texttt{Druhý hráč přesunul objekt na své mřížce.} \\
\\ 
\\ 
\texttt{Jaký je váš další příkaz?} \\
            }
        }
    }
    & \\ \\

    \theutterance \stepcounter{utterance}  
    & & & \multicolumn{4}{p{0.6\linewidth}}{
        \cellcolor[rgb]{0.9,0.9,0.9}{
            \makecell[{{p{\linewidth}}}]{
                \texttt{\tiny{[GM$\langle$P2]}}
                \texttt{POHYB: D, (6,6)} \\
            }
        }
    }
     \\ \\

    \theutterance \stepcounter{utterance}  
    & & & \multicolumn{2}{p{0.3\linewidth}}{
        \cellcolor[rgb]{0.9,0.9,0.9}{
            \makecell[{{p{\linewidth}}}]{
                \texttt{\tiny{[GM$|$GM]}}
                \texttt{POHYB: D, (6,6)} \\
            }
        }
    }
    & & \\ \\

    \theutterance \stepcounter{utterance}  
    & & & \multicolumn{2}{p{0.3\linewidth}}{
        \cellcolor[rgb]{0.9,0.9,0.9}{
            \makecell[{{p{\linewidth}}}]{
                \texttt{\tiny{[GM$|$GM]}}
                \texttt{Penalizace: (6,6) není prázdný, ale obsahuje '┼'.} \\
\\ 
\texttt{Dohromady jste nasbírali 3 z 16 pokut.} \\
\texttt{Zkuste to prosím znovu!} \\
            }
        }
    }
    & & \\ \\

    \theutterance \stepcounter{utterance}  
    & & \multicolumn{4}{p{0.6\linewidth}}{
        \cellcolor[rgb]{0.9,0.9,0.9}{
            \makecell[{{p{\linewidth}}}]{
                \texttt{\tiny{[GM$\rangle$P2]}}
                \texttt{Penalizace: (6,6) není prázdný, ale obsahuje '┼'.} \\
\\ 
\texttt{Dohromady jste nasbírali 3 z 16 pokut.} \\
\texttt{Zkuste to prosím znovu!} \\
            }
        }
    }
    & \\ \\

    \theutterance \stepcounter{utterance}  
    & & & \multicolumn{4}{p{0.6\linewidth}}{
        \cellcolor[rgb]{0.9,0.9,0.9}{
            \makecell[{{p{\linewidth}}}]{
                \texttt{\tiny{[GM$\langle$P2]}}
                \texttt{POHYB: D, (6,7)} \\
            }
        }
    }
     \\ \\

    \theutterance \stepcounter{utterance}  
    & & & \multicolumn{2}{p{0.3\linewidth}}{
        \cellcolor[rgb]{0.9,0.9,0.9}{
            \makecell[{{p{\linewidth}}}]{
                \texttt{\tiny{[GM$|$GM]}}
                \texttt{POHYB: D, (6,7)} \\
            }
        }
    }
    & & \\ \\

    \theutterance \stepcounter{utterance}  
    & & & \multicolumn{2}{p{0.3\linewidth}}{
        \cellcolor[rgb]{0.9,0.9,0.9}{
            \makecell[{{p{\linewidth}}}]{
                \texttt{\tiny{[GM$|$GM]}}
                \texttt{Penalizace: (6,7) není prázdný, ale obsahuje '│'.} \\
\\ 
\texttt{Dohromady jste nasbírali 4 z 16 pokut.} \\
\texttt{Zkuste to prosím znovu!} \\
            }
        }
    }
    & & \\ \\

    \theutterance \stepcounter{utterance}  
    & & \multicolumn{4}{p{0.6\linewidth}}{
        \cellcolor[rgb]{0.9,0.9,0.9}{
            \makecell[{{p{\linewidth}}}]{
                \texttt{\tiny{[GM$\rangle$P2]}}
                \texttt{Penalizace: (6,7) není prázdný, ale obsahuje '│'.} \\
\\ 
\texttt{Dohromady jste nasbírali 4 z 16 pokut.} \\
\texttt{Zkuste to prosím znovu!} \\
            }
        }
    }
    & \\ \\

    \theutterance \stepcounter{utterance}  
    & & & \multicolumn{4}{p{0.6\linewidth}}{
        \cellcolor[rgb]{0.9,0.9,0.9}{
            \makecell[{{p{\linewidth}}}]{
                \texttt{\tiny{[GM$\langle$P2]}}
                \texttt{POHYB: D, (7,6)} \\
            }
        }
    }
     \\ \\

    \theutterance \stepcounter{utterance}  
    & & & \multicolumn{2}{p{0.3\linewidth}}{
        \cellcolor[rgb]{0.9,0.9,0.9}{
            \makecell[{{p{\linewidth}}}]{
                \texttt{\tiny{[GM$|$GM]}}
                \texttt{POHYB: D, (7,6)} \\
            }
        }
    }
    & & \\ \\

    \theutterance \stepcounter{utterance}  
    & & & \multicolumn{2}{p{0.3\linewidth}}{
        \cellcolor[rgb]{0.9,0.9,0.9}{
            \makecell[{{p{\linewidth}}}]{
                \texttt{\tiny{[GM$|$GM]}}
                \texttt{Penalizace: (7,6) není prázdný, ale obsahuje '─'.} \\
\\ 
\texttt{Dohromady jste nasbírali 5 z 16 pokut.} \\
\texttt{Zkuste to prosím znovu!} \\
            }
        }
    }
    & & \\ \\

    \theutterance \stepcounter{utterance}  
    & & \multicolumn{4}{p{0.6\linewidth}}{
        \cellcolor[rgb]{0.9,0.9,0.9}{
            \makecell[{{p{\linewidth}}}]{
                \texttt{\tiny{[GM$\rangle$P2]}}
                \texttt{Penalizace: (7,6) není prázdný, ale obsahuje '─'.} \\
\\ 
\texttt{Dohromady jste nasbírali 5 z 16 pokut.} \\
\texttt{Zkuste to prosím znovu!} \\
            }
        }
    }
    & \\ \\

    \theutterance \stepcounter{utterance}  
    & & & \multicolumn{4}{p{0.6\linewidth}}{
        \cellcolor[rgb]{0.9,0.9,0.9}{
            \makecell[{{p{\linewidth}}}]{
                \texttt{\tiny{[GM$\langle$P2]}}
                \texttt{POHYB: D, (7,7)} \\
            }
        }
    }
     \\ \\

    \theutterance \stepcounter{utterance}  
    & & & \multicolumn{2}{p{0.3\linewidth}}{
        \cellcolor[rgb]{0.9,0.9,0.9}{
            \makecell[{{p{\linewidth}}}]{
                \texttt{\tiny{[GM$|$GM]}}
                \texttt{POHYB: D, (7,7)} \\
            }
        }
    }
    & & \\ \\

    \theutterance \stepcounter{utterance}  
    & & & \multicolumn{2}{p{0.3\linewidth}}{
        \cellcolor[rgb]{0.9,0.9,0.9}{
            \makecell[{{p{\linewidth}}}]{
                \texttt{\tiny{[GM$|$GM]}}
                \texttt{D moved to (7, 7)} \\
            }
        }
    }
    & & \\ \\

    \theutterance \stepcounter{utterance}  
    & & \multicolumn{4}{p{0.6\linewidth}}{
        \cellcolor[rgb]{0.9,0.9,0.9}{
            \makecell[{{p{\linewidth}}}]{
                \texttt{\tiny{[P1$\langle$GM]}}
                \texttt{Úspěšně přesunuto 'P' do (1,1). Vaše aktualizovaná mřížka vypadá takto:} \\
\texttt{```} \\
\texttt{ 1234567} \\
\texttt{╔═╤═════╗} \\
\texttt{║P│◌A◌◌◌║ 1} \\
\texttt{╟─┤◌T◌◌◌║ 2} \\
\texttt{║O└─┐◌◌◌║ 3} \\
\texttt{║◌◌◌├─┐◌║ 4} \\
\texttt{║◌◌S│◌│◌║ 5} \\
\texttt{╟─┬─┴─┼─╢ 6} \\
\texttt{║◌│◌DM│◌║ 7} \\
\texttt{╚═╧═══╧═╝} \\
\texttt{```} \\
\texttt{Právě hrajete kolo 4 z maximálně 28.} \\
\texttt{Dohromady jste nasbírali 5 z 16 pokut.} \\
\texttt{Druhý hráč přesunul objekt na své mřížce.} \\
\\ 
\\ 
\texttt{Jaký je váš další příkaz?} \\
            }
        }
    }
    & \\ \\

    \theutterance \stepcounter{utterance}  
    & \multicolumn{4}{p{0.6\linewidth}}{
        \cellcolor[rgb]{0.9,0.9,0.9}{
            \makecell[{{p{\linewidth}}}]{
                \texttt{\tiny{[P1$\rangle$GM]}}
                \texttt{POHYB: A, (7,1)} \\
            }
        }
    }
    & & \\ \\

    \theutterance \stepcounter{utterance}  
    & & & \multicolumn{2}{p{0.3\linewidth}}{
        \cellcolor[rgb]{0.9,0.9,0.9}{
            \makecell[{{p{\linewidth}}}]{
                \texttt{\tiny{[GM$|$GM]}}
                \texttt{POHYB: A, (7,1)} \\
            }
        }
    }
    & & \\ \\

    \theutterance \stepcounter{utterance}  
    & & & \multicolumn{2}{p{0.3\linewidth}}{
        \cellcolor[rgb]{0.9,0.9,0.9}{
            \makecell[{{p{\linewidth}}}]{
                \texttt{\tiny{[GM$|$GM]}}
                \texttt{A moved to (7, 1)} \\
            }
        }
    }
    & & \\ \\

    \theutterance \stepcounter{utterance}  
    & & \multicolumn{4}{p{0.6\linewidth}}{
        \cellcolor[rgb]{0.9,0.9,0.9}{
            \makecell[{{p{\linewidth}}}]{
                \texttt{\tiny{[GM$\rangle$P2]}}
                \texttt{Úspěšně přesunuto 'D' do (7,7). Vaše aktualizovaná mřížka vypadá takto:} \\
\texttt{```} \\
\texttt{ 1234567} \\
\texttt{╔═╤═════╗} \\
\texttt{║P│◌◌◌◌◌║ 1} \\
\texttt{╟─┤◌◌◌◌◌║ 2} \\
\texttt{║◌└─┐◌A◌║ 3} \\
\texttt{║◌◌M├─┐◌║ 4} \\
\texttt{║◌T◌│S│O║ 5} \\
\texttt{╟─┬─┴─┼─╢ 6} \\
\texttt{║◌│◌◌◌│D║ 7} \\
\texttt{╚═╧═══╧═╝} \\
\texttt{```} \\
\texttt{Právě hrajete kolo 5 z maximálně 28.} \\
\texttt{Dohromady jste nasbírali 5 z 16 pokut.} \\
\texttt{Druhý hráč přesunul objekt na své mřížce.} \\
\\ 
\\ 
\texttt{Jaký je váš další příkaz?} \\
            }
        }
    }
    & \\ \\

    \theutterance \stepcounter{utterance}  
    & & & \multicolumn{4}{p{0.6\linewidth}}{
        \cellcolor[rgb]{0.9,0.9,0.9}{
            \makecell[{{p{\linewidth}}}]{
                \texttt{\tiny{[GM$\langle$P2]}}
                \texttt{POHYB: A, (7,1)} \\
            }
        }
    }
     \\ \\

    \theutterance \stepcounter{utterance}  
    & & & \multicolumn{2}{p{0.3\linewidth}}{
        \cellcolor[rgb]{0.9,0.9,0.9}{
            \makecell[{{p{\linewidth}}}]{
                \texttt{\tiny{[GM$|$GM]}}
                \texttt{POHYB: A, (7,1)} \\
            }
        }
    }
    & & \\ \\

    \theutterance \stepcounter{utterance}  
    & & & \multicolumn{2}{p{0.3\linewidth}}{
        \cellcolor[rgb]{0.9,0.9,0.9}{
            \makecell[{{p{\linewidth}}}]{
                \texttt{\tiny{[GM$|$GM]}}
                \texttt{A moved to (7, 1)} \\
            }
        }
    }
    & & \\ \\

    \theutterance \stepcounter{utterance}  
    & & \multicolumn{4}{p{0.6\linewidth}}{
        \cellcolor[rgb]{0.9,0.9,0.9}{
            \makecell[{{p{\linewidth}}}]{
                \texttt{\tiny{[P1$\langle$GM]}}
                \texttt{Úspěšně přesunuto 'A' do (7,1). Vaše aktualizovaná mřížka vypadá takto:} \\
\texttt{```} \\
\texttt{ 1234567} \\
\texttt{╔═╤═════╗} \\
\texttt{║P│◌◌◌◌A║ 1} \\
\texttt{╟─┤◌T◌◌◌║ 2} \\
\texttt{║O└─┐◌◌◌║ 3} \\
\texttt{║◌◌◌├─┐◌║ 4} \\
\texttt{║◌◌S│◌│◌║ 5} \\
\texttt{╟─┬─┴─┼─╢ 6} \\
\texttt{║◌│◌DM│◌║ 7} \\
\texttt{╚═╧═══╧═╝} \\
\texttt{```} \\
\texttt{Právě hrajete kolo 5 z maximálně 28.} \\
\texttt{Dohromady jste nasbírali 5 z 16 pokut.} \\
\texttt{Druhý hráč přesunul objekt na své mřížce.} \\
\\ 
\\ 
\texttt{Jaký je váš další příkaz?} \\
            }
        }
    }
    & \\ \\

    \theutterance \stepcounter{utterance}  
    & \multicolumn{4}{p{0.6\linewidth}}{
        \cellcolor[rgb]{0.9,0.9,0.9}{
            \makecell[{{p{\linewidth}}}]{
                \texttt{\tiny{[P1$\rangle$GM]}}
                \texttt{POHYB: T, (3,3)} \\
            }
        }
    }
    & & \\ \\

    \theutterance \stepcounter{utterance}  
    & & & \multicolumn{2}{p{0.3\linewidth}}{
        \cellcolor[rgb]{0.9,0.9,0.9}{
            \makecell[{{p{\linewidth}}}]{
                \texttt{\tiny{[GM$|$GM]}}
                \texttt{POHYB: T, (3,3)} \\
            }
        }
    }
    & & \\ \\

    \theutterance \stepcounter{utterance}  
    & & & \multicolumn{2}{p{0.3\linewidth}}{
        \cellcolor[rgb]{0.9,0.9,0.9}{
            \makecell[{{p{\linewidth}}}]{
                \texttt{\tiny{[GM$|$GM]}}
                \texttt{Penalizace: (3,3) není prázdný, ale obsahuje '─'.} \\
\\ 
\texttt{Dohromady jste nasbírali 6 z 16 pokut.} \\
\texttt{Zkuste to prosím znovu!} \\
            }
        }
    }
    & & \\ \\

    \theutterance \stepcounter{utterance}  
    & & \multicolumn{4}{p{0.6\linewidth}}{
        \cellcolor[rgb]{0.9,0.9,0.9}{
            \makecell[{{p{\linewidth}}}]{
                \texttt{\tiny{[P1$\langle$GM]}}
                \texttt{Penalizace: (3,3) není prázdný, ale obsahuje '─'.} \\
\\ 
\texttt{Dohromady jste nasbírali 6 z 16 pokut.} \\
\texttt{Zkuste to prosím znovu!} \\
            }
        }
    }
    & \\ \\

    \theutterance \stepcounter{utterance}  
    & \multicolumn{4}{p{0.6\linewidth}}{
        \cellcolor[rgb]{0.9,0.9,0.9}{
            \makecell[{{p{\linewidth}}}]{
                \texttt{\tiny{[P1$\rangle$GM]}}
                \texttt{POHYB: T, (2,3)} \\
            }
        }
    }
    & & \\ \\

    \theutterance \stepcounter{utterance}  
    & & & \multicolumn{2}{p{0.3\linewidth}}{
        \cellcolor[rgb]{0.9,0.9,0.9}{
            \makecell[{{p{\linewidth}}}]{
                \texttt{\tiny{[GM$|$GM]}}
                \texttt{POHYB: T, (2,3)} \\
            }
        }
    }
    & & \\ \\

    \theutterance \stepcounter{utterance}  
    & & & \multicolumn{2}{p{0.3\linewidth}}{
        \cellcolor[rgb]{0.9,0.9,0.9}{
            \makecell[{{p{\linewidth}}}]{
                \texttt{\tiny{[GM$|$GM]}}
                \texttt{Penalizace: (2,3) není prázdný, ale obsahuje '└'.} \\
\\ 
\texttt{Dohromady jste nasbírali 7 z 16 pokut.} \\
\texttt{Zkuste to prosím znovu!} \\
            }
        }
    }
    & & \\ \\

    \theutterance \stepcounter{utterance}  
    & & \multicolumn{4}{p{0.6\linewidth}}{
        \cellcolor[rgb]{0.9,0.9,0.9}{
            \makecell[{{p{\linewidth}}}]{
                \texttt{\tiny{[P1$\langle$GM]}}
                \texttt{Penalizace: (2,3) není prázdný, ale obsahuje '└'.} \\
\\ 
\texttt{Dohromady jste nasbírali 7 z 16 pokut.} \\
\texttt{Zkuste to prosím znovu!} \\
            }
        }
    }
    & \\ \\

    \theutterance \stepcounter{utterance}  
    & \multicolumn{4}{p{0.6\linewidth}}{
        \cellcolor[rgb]{0.9,0.9,0.9}{
            \makecell[{{p{\linewidth}}}]{
                \texttt{\tiny{[P1$\rangle$GM]}}
                \texttt{POHYB: T, (4,3)} \\
            }
        }
    }
    & & \\ \\

    \theutterance \stepcounter{utterance}  
    & & & \multicolumn{2}{p{0.3\linewidth}}{
        \cellcolor[rgb]{0.9,0.9,0.9}{
            \makecell[{{p{\linewidth}}}]{
                \texttt{\tiny{[GM$|$GM]}}
                \texttt{POHYB: T, (4,3)} \\
            }
        }
    }
    & & \\ \\

    \theutterance \stepcounter{utterance}  
    & & & \multicolumn{2}{p{0.3\linewidth}}{
        \cellcolor[rgb]{0.9,0.9,0.9}{
            \makecell[{{p{\linewidth}}}]{
                \texttt{\tiny{[GM$|$GM]}}
                \texttt{Penalizace: (4,3) není prázdný, ale obsahuje '┐'.} \\
\\ 
\texttt{Dohromady jste nasbírali 8 z 16 pokut.} \\
\texttt{Zkuste to prosím znovu!} \\
            }
        }
    }
    & & \\ \\

    \theutterance \stepcounter{utterance}  
    & & \multicolumn{4}{p{0.6\linewidth}}{
        \cellcolor[rgb]{0.9,0.9,0.9}{
            \makecell[{{p{\linewidth}}}]{
                \texttt{\tiny{[P1$\langle$GM]}}
                \texttt{Penalizace: (4,3) není prázdný, ale obsahuje '┐'.} \\
\\ 
\texttt{Dohromady jste nasbírali 8 z 16 pokut.} \\
\texttt{Zkuste to prosím znovu!} \\
            }
        }
    }
    & \\ \\

    \theutterance \stepcounter{utterance}  
    & \multicolumn{4}{p{0.6\linewidth}}{
        \cellcolor[rgb]{0.9,0.9,0.9}{
            \makecell[{{p{\linewidth}}}]{
                \texttt{\tiny{[P1$\rangle$GM]}}
                \texttt{POHYB: T, (5,3)} \\
            }
        }
    }
    & & \\ \\

    \theutterance \stepcounter{utterance}  
    & & & \multicolumn{2}{p{0.3\linewidth}}{
        \cellcolor[rgb]{0.9,0.9,0.9}{
            \makecell[{{p{\linewidth}}}]{
                \texttt{\tiny{[GM$|$GM]}}
                \texttt{POHYB: T, (5,3)} \\
            }
        }
    }
    & & \\ \\

    \theutterance \stepcounter{utterance}  
    & & & \multicolumn{2}{p{0.3\linewidth}}{
        \cellcolor[rgb]{0.9,0.9,0.9}{
            \makecell[{{p{\linewidth}}}]{
                \texttt{\tiny{[GM$|$GM]}}
                \texttt{T moved to (5, 3)} \\
            }
        }
    }
    & & \\ \\

    \theutterance \stepcounter{utterance}  
    & & \multicolumn{4}{p{0.6\linewidth}}{
        \cellcolor[rgb]{0.9,0.9,0.9}{
            \makecell[{{p{\linewidth}}}]{
                \texttt{\tiny{[GM$\rangle$P2]}}
                \texttt{Úspěšně přesunuto 'A' do (7,1). Vaše aktualizovaná mřížka vypadá takto:} \\
\texttt{```} \\
\texttt{ 1234567} \\
\texttt{╔═╤═════╗} \\
\texttt{║P│◌◌◌◌A║ 1} \\
\texttt{╟─┤◌◌◌◌◌║ 2} \\
\texttt{║◌└─┐◌◌◌║ 3} \\
\texttt{║◌◌M├─┐◌║ 4} \\
\texttt{║◌T◌│S│O║ 5} \\
\texttt{╟─┬─┴─┼─╢ 6} \\
\texttt{║◌│◌◌◌│D║ 7} \\
\texttt{╚═╧═══╧═╝} \\
\texttt{```} \\
\texttt{Právě hrajete kolo 6 z maximálně 28.} \\
\texttt{Dohromady jste nasbírali 8 z 16 pokut.} \\
\texttt{Druhý hráč přesunul objekt na své mřížce.} \\
\\ 
\\ 
\texttt{Jaký je váš další příkaz?} \\
            }
        }
    }
    & \\ \\

    \theutterance \stepcounter{utterance}  
    & & & \multicolumn{4}{p{0.6\linewidth}}{
        \cellcolor[rgb]{0.9,0.9,0.9}{
            \makecell[{{p{\linewidth}}}]{
                \texttt{\tiny{[GM$\langle$P2]}}
                \texttt{POHYB: O, (1,3)} \\
            }
        }
    }
     \\ \\

    \theutterance \stepcounter{utterance}  
    & & & \multicolumn{2}{p{0.3\linewidth}}{
        \cellcolor[rgb]{0.9,0.9,0.9}{
            \makecell[{{p{\linewidth}}}]{
                \texttt{\tiny{[GM$|$GM]}}
                \texttt{POHYB: O, (1,3)} \\
            }
        }
    }
    & & \\ \\

    \theutterance \stepcounter{utterance}  
    & & & \multicolumn{2}{p{0.3\linewidth}}{
        \cellcolor[rgb]{0.9,0.9,0.9}{
            \makecell[{{p{\linewidth}}}]{
                \texttt{\tiny{[GM$|$GM]}}
                \texttt{O moved to (1, 3)} \\
            }
        }
    }
    & & \\ \\

    \theutterance \stepcounter{utterance}  
    & & \multicolumn{4}{p{0.6\linewidth}}{
        \cellcolor[rgb]{0.9,0.9,0.9}{
            \makecell[{{p{\linewidth}}}]{
                \texttt{\tiny{[P1$\langle$GM]}}
                \texttt{Úspěšně přesunuto 'T' do (5,3). Vaše aktualizovaná mřížka vypadá takto:} \\
\texttt{```} \\
\texttt{ 1234567} \\
\texttt{╔═╤═════╗} \\
\texttt{║P│◌◌◌◌A║ 1} \\
\texttt{╟─┤◌◌◌◌◌║ 2} \\
\texttt{║O└─┐T◌◌║ 3} \\
\texttt{║◌◌◌├─┐◌║ 4} \\
\texttt{║◌◌S│◌│◌║ 5} \\
\texttt{╟─┬─┴─┼─╢ 6} \\
\texttt{║◌│◌DM│◌║ 7} \\
\texttt{╚═╧═══╧═╝} \\
\texttt{```} \\
\texttt{Právě hrajete kolo 6 z maximálně 28.} \\
\texttt{Dohromady jste nasbírali 8 z 16 pokut.} \\
\texttt{Druhý hráč přesunul objekt na své mřížce.} \\
\\ 
\\ 
\texttt{Jaký je váš další příkaz?} \\
            }
        }
    }
    & \\ \\

    \theutterance \stepcounter{utterance}  
    & \multicolumn{4}{p{0.6\linewidth}}{
        \cellcolor[rgb]{0.9,0.9,0.9}{
            \makecell[{{p{\linewidth}}}]{
                \texttt{\tiny{[P1$\rangle$GM]}}
                \texttt{POHYB: T, (5,2)} \\
            }
        }
    }
    & & \\ \\

    \theutterance \stepcounter{utterance}  
    & & & \multicolumn{2}{p{0.3\linewidth}}{
        \cellcolor[rgb]{0.9,0.9,0.9}{
            \makecell[{{p{\linewidth}}}]{
                \texttt{\tiny{[GM$|$GM]}}
                \texttt{POHYB: T, (5,2)} \\
            }
        }
    }
    & & \\ \\

    \theutterance \stepcounter{utterance}  
    & & & \multicolumn{2}{p{0.3\linewidth}}{
        \cellcolor[rgb]{0.9,0.9,0.9}{
            \makecell[{{p{\linewidth}}}]{
                \texttt{\tiny{[GM$|$GM]}}
                \texttt{T moved to (5, 2)} \\
            }
        }
    }
    & & \\ \\

    \theutterance \stepcounter{utterance}  
    & & \multicolumn{4}{p{0.6\linewidth}}{
        \cellcolor[rgb]{0.9,0.9,0.9}{
            \makecell[{{p{\linewidth}}}]{
                \texttt{\tiny{[GM$\rangle$P2]}}
                \texttt{Úspěšně přesunuto 'O' do (1,3). Vaše aktualizovaná mřížka vypadá takto:} \\
\texttt{```} \\
\texttt{ 1234567} \\
\texttt{╔═╤═════╗} \\
\texttt{║P│◌◌◌◌A║ 1} \\
\texttt{╟─┤◌◌◌◌◌║ 2} \\
\texttt{║O└─┐◌◌◌║ 3} \\
\texttt{║◌◌M├─┐◌║ 4} \\
\texttt{║◌T◌│S│◌║ 5} \\
\texttt{╟─┬─┴─┼─╢ 6} \\
\texttt{║◌│◌◌◌│D║ 7} \\
\texttt{╚═╧═══╧═╝} \\
\texttt{```} \\
\texttt{Právě hrajete kolo 7 z maximálně 28.} \\
\texttt{Dohromady jste nasbírali 8 z 16 pokut.} \\
\texttt{Druhý hráč přesunul objekt na své mřížce.} \\
\\ 
\\ 
\texttt{Jaký je váš další příkaz?} \\
            }
        }
    }
    & \\ \\

    \theutterance \stepcounter{utterance}  
    & & & \multicolumn{4}{p{0.6\linewidth}}{
        \cellcolor[rgb]{0.9,0.9,0.9}{
            \makecell[{{p{\linewidth}}}]{
                \texttt{\tiny{[GM$\langle$P2]}}
                \texttt{POHYB: T, (3,3)} \\
            }
        }
    }
     \\ \\

    \theutterance \stepcounter{utterance}  
    & & & \multicolumn{2}{p{0.3\linewidth}}{
        \cellcolor[rgb]{0.9,0.9,0.9}{
            \makecell[{{p{\linewidth}}}]{
                \texttt{\tiny{[GM$|$GM]}}
                \texttt{POHYB: T, (3,3)} \\
            }
        }
    }
    & & \\ \\

    \theutterance \stepcounter{utterance}  
    & & & \multicolumn{2}{p{0.3\linewidth}}{
        \cellcolor[rgb]{0.9,0.9,0.9}{
            \makecell[{{p{\linewidth}}}]{
                \texttt{\tiny{[GM$|$GM]}}
                \texttt{Penalizace: (3,3) není prázdný, ale obsahuje '─'.} \\
\\ 
\texttt{Dohromady jste nasbírali 9 z 16 pokut.} \\
\texttt{Zkuste to prosím znovu!} \\
            }
        }
    }
    & & \\ \\

    \theutterance \stepcounter{utterance}  
    & & \multicolumn{4}{p{0.6\linewidth}}{
        \cellcolor[rgb]{0.9,0.9,0.9}{
            \makecell[{{p{\linewidth}}}]{
                \texttt{\tiny{[GM$\rangle$P2]}}
                \texttt{Penalizace: (3,3) není prázdný, ale obsahuje '─'.} \\
\\ 
\texttt{Dohromady jste nasbírali 9 z 16 pokut.} \\
\texttt{Zkuste to prosím znovu!} \\
            }
        }
    }
    & \\ \\

    \theutterance \stepcounter{utterance}  
    & & & \multicolumn{4}{p{0.6\linewidth}}{
        \cellcolor[rgb]{0.9,0.9,0.9}{
            \makecell[{{p{\linewidth}}}]{
                \texttt{\tiny{[GM$\langle$P2]}}
                \texttt{POHYB: T, (3,2)} \\
            }
        }
    }
     \\ \\

    \theutterance \stepcounter{utterance}  
    & & & \multicolumn{2}{p{0.3\linewidth}}{
        \cellcolor[rgb]{0.9,0.9,0.9}{
            \makecell[{{p{\linewidth}}}]{
                \texttt{\tiny{[GM$|$GM]}}
                \texttt{POHYB: T, (3,2)} \\
            }
        }
    }
    & & \\ \\

    \theutterance \stepcounter{utterance}  
    & & & \multicolumn{2}{p{0.3\linewidth}}{
        \cellcolor[rgb]{0.9,0.9,0.9}{
            \makecell[{{p{\linewidth}}}]{
                \texttt{\tiny{[GM$|$GM]}}
                \texttt{T moved to (3, 2)} \\
            }
        }
    }
    & & \\ \\

    \theutterance \stepcounter{utterance}  
    & & \multicolumn{4}{p{0.6\linewidth}}{
        \cellcolor[rgb]{0.9,0.9,0.9}{
            \makecell[{{p{\linewidth}}}]{
                \texttt{\tiny{[P1$\langle$GM]}}
                \texttt{Úspěšně přesunuto 'T' do (5,2). Vaše aktualizovaná mřížka vypadá takto:} \\
\texttt{```} \\
\texttt{ 1234567} \\
\texttt{╔═╤═════╗} \\
\texttt{║P│◌◌◌◌A║ 1} \\
\texttt{╟─┤◌◌T◌◌║ 2} \\
\texttt{║O└─┐◌◌◌║ 3} \\
\texttt{║◌◌◌├─┐◌║ 4} \\
\texttt{║◌◌S│◌│◌║ 5} \\
\texttt{╟─┬─┴─┼─╢ 6} \\
\texttt{║◌│◌DM│◌║ 7} \\
\texttt{╚═╧═══╧═╝} \\
\texttt{```} \\
\texttt{Právě hrajete kolo 7 z maximálně 28.} \\
\texttt{Dohromady jste nasbírali 9 z 16 pokut.} \\
\texttt{Druhý hráč přesunul objekt na své mřížce.} \\
\\ 
\\ 
\texttt{Jaký je váš další příkaz?} \\
            }
        }
    }
    & \\ \\

    \theutterance \stepcounter{utterance}  
    & \multicolumn{4}{p{0.6\linewidth}}{
        \cellcolor[rgb]{0.9,0.9,0.9}{
            \makecell[{{p{\linewidth}}}]{
                \texttt{\tiny{[P1$\rangle$GM]}}
                \texttt{POHYB: O, (1,3)} \\
            }
        }
    }
    & & \\ \\

    \theutterance \stepcounter{utterance}  
    & & & \multicolumn{2}{p{0.3\linewidth}}{
        \cellcolor[rgb]{0.9,0.9,0.9}{
            \makecell[{{p{\linewidth}}}]{
                \texttt{\tiny{[GM$|$GM]}}
                \texttt{POHYB: O, (1,3)} \\
            }
        }
    }
    & & \\ \\

    \theutterance \stepcounter{utterance}  
    & & & \multicolumn{2}{p{0.3\linewidth}}{
        \cellcolor[rgb]{0.9,0.9,0.9}{
            \makecell[{{p{\linewidth}}}]{
                \texttt{\tiny{[GM$|$GM]}}
                \texttt{Penalizace: (1,3) není prázdný, ale obsahuje 'O'.} \\
\\ 
\texttt{Dohromady jste nasbírali 10 z 16 pokut.} \\
\texttt{Zkuste to prosím znovu!} \\
            }
        }
    }
    & & \\ \\

    \theutterance \stepcounter{utterance}  
    & & \multicolumn{4}{p{0.6\linewidth}}{
        \cellcolor[rgb]{0.9,0.9,0.9}{
            \makecell[{{p{\linewidth}}}]{
                \texttt{\tiny{[P1$\langle$GM]}}
                \texttt{Penalizace: (1,3) není prázdný, ale obsahuje 'O'.} \\
\\ 
\texttt{Dohromady jste nasbírali 10 z 16 pokut.} \\
\texttt{Zkuste to prosím znovu!} \\
            }
        }
    }
    & \\ \\

    \theutterance \stepcounter{utterance}  
    & \multicolumn{4}{p{0.6\linewidth}}{
        \cellcolor[rgb]{0.9,0.9,0.9}{
            \makecell[{{p{\linewidth}}}]{
                \texttt{\tiny{[P1$\rangle$GM]}}
                \texttt{POHYB: S, (5,5)} \\
            }
        }
    }
    & & \\ \\

    \theutterance \stepcounter{utterance}  
    & & & \multicolumn{2}{p{0.3\linewidth}}{
        \cellcolor[rgb]{0.9,0.9,0.9}{
            \makecell[{{p{\linewidth}}}]{
                \texttt{\tiny{[GM$|$GM]}}
                \texttt{POHYB: S, (5,5)} \\
            }
        }
    }
    & & \\ \\

    \theutterance \stepcounter{utterance}  
    & & & \multicolumn{2}{p{0.3\linewidth}}{
        \cellcolor[rgb]{0.9,0.9,0.9}{
            \makecell[{{p{\linewidth}}}]{
                \texttt{\tiny{[GM$|$GM]}}
                \texttt{S moved to (5, 5)} \\
            }
        }
    }
    & & \\ \\

    \theutterance \stepcounter{utterance}  
    & & \multicolumn{4}{p{0.6\linewidth}}{
        \cellcolor[rgb]{0.9,0.9,0.9}{
            \makecell[{{p{\linewidth}}}]{
                \texttt{\tiny{[GM$\rangle$P2]}}
                \texttt{Úspěšně přesunuto 'T' do (3,2). Vaše aktualizovaná mřížka vypadá takto:} \\
\texttt{```} \\
\texttt{ 1234567} \\
\texttt{╔═╤═════╗} \\
\texttt{║P│◌◌◌◌A║ 1} \\
\texttt{╟─┤T◌◌◌◌║ 2} \\
\texttt{║O└─┐◌◌◌║ 3} \\
\texttt{║◌◌M├─┐◌║ 4} \\
\texttt{║◌◌◌│S│◌║ 5} \\
\texttt{╟─┬─┴─┼─╢ 6} \\
\texttt{║◌│◌◌◌│D║ 7} \\
\texttt{╚═╧═══╧═╝} \\
\texttt{```} \\
\texttt{Právě hrajete kolo 8 z maximálně 28.} \\
\texttt{Dohromady jste nasbírali 10 z 16 pokut.} \\
\texttt{Druhý hráč přesunul objekt na své mřížce.} \\
\\ 
\\ 
\texttt{Jaký je váš další příkaz?} \\
            }
        }
    }
    & \\ \\

    \theutterance \stepcounter{utterance}  
    & & & \multicolumn{4}{p{0.6\linewidth}}{
        \cellcolor[rgb]{0.9,0.9,0.9}{
            \makecell[{{p{\linewidth}}}]{
                \texttt{\tiny{[GM$\langle$P2]}}
                \texttt{POHYB: M, (5,7)} \\
            }
        }
    }
     \\ \\

    \theutterance \stepcounter{utterance}  
    & & & \multicolumn{2}{p{0.3\linewidth}}{
        \cellcolor[rgb]{0.9,0.9,0.9}{
            \makecell[{{p{\linewidth}}}]{
                \texttt{\tiny{[GM$|$GM]}}
                \texttt{POHYB: M, (5,7)} \\
            }
        }
    }
    & & \\ \\

    \theutterance \stepcounter{utterance}  
    & & & \multicolumn{2}{p{0.3\linewidth}}{
        \cellcolor[rgb]{0.9,0.9,0.9}{
            \makecell[{{p{\linewidth}}}]{
                \texttt{\tiny{[GM$|$GM]}}
                \texttt{M moved to (5, 7)} \\
            }
        }
    }
    & & \\ \\

    \theutterance \stepcounter{utterance}  
    & & \multicolumn{4}{p{0.6\linewidth}}{
        \cellcolor[rgb]{0.9,0.9,0.9}{
            \makecell[{{p{\linewidth}}}]{
                \texttt{\tiny{[P1$\langle$GM]}}
                \texttt{Úspěšně přesunuto 'S' do (5,5). Vaše aktualizovaná mřížka vypadá takto:} \\
\texttt{```} \\
\texttt{ 1234567} \\
\texttt{╔═╤═════╗} \\
\texttt{║P│◌◌◌◌A║ 1} \\
\texttt{╟─┤◌◌T◌◌║ 2} \\
\texttt{║O└─┐◌◌◌║ 3} \\
\texttt{║◌◌◌├─┐◌║ 4} \\
\texttt{║◌◌◌│S│◌║ 5} \\
\texttt{╟─┬─┴─┼─╢ 6} \\
\texttt{║◌│◌DM│◌║ 7} \\
\texttt{╚═╧═══╧═╝} \\
\texttt{```} \\
\texttt{Právě hrajete kolo 8 z maximálně 28.} \\
\texttt{Dohromady jste nasbírali 10 z 16 pokut.} \\
\texttt{Druhý hráč přesunul objekt na své mřížce.} \\
\\ 
\\ 
\texttt{Jaký je váš další příkaz?} \\
            }
        }
    }
    & \\ \\

    \theutterance \stepcounter{utterance}  
    & \multicolumn{4}{p{0.6\linewidth}}{
        \cellcolor[rgb]{0.9,0.9,0.9}{
            \makecell[{{p{\linewidth}}}]{
                \texttt{\tiny{[P1$\rangle$GM]}}
                \texttt{POHYB: D, (6,6)} \\
            }
        }
    }
    & & \\ \\

    \theutterance \stepcounter{utterance}  
    & & & \multicolumn{2}{p{0.3\linewidth}}{
        \cellcolor[rgb]{0.9,0.9,0.9}{
            \makecell[{{p{\linewidth}}}]{
                \texttt{\tiny{[GM$|$GM]}}
                \texttt{POHYB: D, (6,6)} \\
            }
        }
    }
    & & \\ \\

    \theutterance \stepcounter{utterance}  
    & & & \multicolumn{2}{p{0.3\linewidth}}{
        \cellcolor[rgb]{0.9,0.9,0.9}{
            \makecell[{{p{\linewidth}}}]{
                \texttt{\tiny{[GM$|$GM]}}
                \texttt{Penalizace: (6,6) není prázdný, ale obsahuje '┼'.} \\
\\ 
\texttt{Dohromady jste nasbírali 11 z 16 pokut.} \\
\texttt{Zkuste to prosím znovu!} \\
            }
        }
    }
    & & \\ \\

    \theutterance \stepcounter{utterance}  
    & & \multicolumn{4}{p{0.6\linewidth}}{
        \cellcolor[rgb]{0.9,0.9,0.9}{
            \makecell[{{p{\linewidth}}}]{
                \texttt{\tiny{[P1$\langle$GM]}}
                \texttt{Penalizace: (6,6) není prázdný, ale obsahuje '┼'.} \\
\\ 
\texttt{Dohromady jste nasbírali 11 z 16 pokut.} \\
\texttt{Zkuste to prosím znovu!} \\
            }
        }
    }
    & \\ \\

    \theutterance \stepcounter{utterance}  
    & \multicolumn{4}{p{0.6\linewidth}}{
        \cellcolor[rgb]{0.9,0.9,0.9}{
            \makecell[{{p{\linewidth}}}]{
                \texttt{\tiny{[P1$\rangle$GM]}}
                \texttt{POHYB: D, (6,4)} \\
            }
        }
    }
    & & \\ \\

    \theutterance \stepcounter{utterance}  
    & & & \multicolumn{2}{p{0.3\linewidth}}{
        \cellcolor[rgb]{0.9,0.9,0.9}{
            \makecell[{{p{\linewidth}}}]{
                \texttt{\tiny{[GM$|$GM]}}
                \texttt{POHYB: D, (6,4)} \\
            }
        }
    }
    & & \\ \\

    \theutterance \stepcounter{utterance}  
    & & & \multicolumn{2}{p{0.3\linewidth}}{
        \cellcolor[rgb]{0.9,0.9,0.9}{
            \makecell[{{p{\linewidth}}}]{
                \texttt{\tiny{[GM$|$GM]}}
                \texttt{Penalizace: (6,4) není prázdný, ale obsahuje '┐'.} \\
\\ 
\texttt{Dohromady jste nasbírali 12 z 16 pokut.} \\
\texttt{Zkuste to prosím znovu!} \\
            }
        }
    }
    & & \\ \\

    \theutterance \stepcounter{utterance}  
    & & \multicolumn{4}{p{0.6\linewidth}}{
        \cellcolor[rgb]{0.9,0.9,0.9}{
            \makecell[{{p{\linewidth}}}]{
                \texttt{\tiny{[P1$\langle$GM]}}
                \texttt{Penalizace: (6,4) není prázdný, ale obsahuje '┐'.} \\
\\ 
\texttt{Dohromady jste nasbírali 12 z 16 pokut.} \\
\texttt{Zkuste to prosím znovu!} \\
            }
        }
    }
    & \\ \\

    \theutterance \stepcounter{utterance}  
    & \multicolumn{4}{p{0.6\linewidth}}{
        \cellcolor[rgb]{0.9,0.9,0.9}{
            \makecell[{{p{\linewidth}}}]{
                \texttt{\tiny{[P1$\rangle$GM]}}
                \texttt{POHYB: D, (5,6)} \\
            }
        }
    }
    & & \\ \\

    \theutterance \stepcounter{utterance}  
    & & & \multicolumn{2}{p{0.3\linewidth}}{
        \cellcolor[rgb]{0.9,0.9,0.9}{
            \makecell[{{p{\linewidth}}}]{
                \texttt{\tiny{[GM$|$GM]}}
                \texttt{POHYB: D, (5,6)} \\
            }
        }
    }
    & & \\ \\

    \theutterance \stepcounter{utterance}  
    & & & \multicolumn{2}{p{0.3\linewidth}}{
        \cellcolor[rgb]{0.9,0.9,0.9}{
            \makecell[{{p{\linewidth}}}]{
                \texttt{\tiny{[GM$|$GM]}}
                \texttt{Penalizace: (5,6) není prázdný, ale obsahuje '─'.} \\
\\ 
\texttt{Dohromady jste nasbírali 13 z 16 pokut.} \\
\texttt{Zkuste to prosím znovu!} \\
            }
        }
    }
    & & \\ \\

    \theutterance \stepcounter{utterance}  
    & & \multicolumn{4}{p{0.6\linewidth}}{
        \cellcolor[rgb]{0.9,0.9,0.9}{
            \makecell[{{p{\linewidth}}}]{
                \texttt{\tiny{[P1$\langle$GM]}}
                \texttt{Penalizace: (5,6) není prázdný, ale obsahuje '─'.} \\
\\ 
\texttt{Dohromady jste nasbírali 13 z 16 pokut.} \\
\texttt{Zkuste to prosím znovu!} \\
            }
        }
    }
    & \\ \\

    \theutterance \stepcounter{utterance}  
    & \multicolumn{4}{p{0.6\linewidth}}{
        \cellcolor[rgb]{0.9,0.9,0.9}{
            \makecell[{{p{\linewidth}}}]{
                \texttt{\tiny{[P1$\rangle$GM]}}
                \texttt{POHYB: D, (7,6)} \\
            }
        }
    }
    & & \\ \\

    \theutterance \stepcounter{utterance}  
    & & & \multicolumn{2}{p{0.3\linewidth}}{
        \cellcolor[rgb]{0.9,0.9,0.9}{
            \makecell[{{p{\linewidth}}}]{
                \texttt{\tiny{[GM$|$GM]}}
                \texttt{POHYB: D, (7,6)} \\
            }
        }
    }
    & & \\ \\

    \theutterance \stepcounter{utterance}  
    & & & \multicolumn{2}{p{0.3\linewidth}}{
        \cellcolor[rgb]{0.9,0.9,0.9}{
            \makecell[{{p{\linewidth}}}]{
                \texttt{\tiny{[GM$|$GM]}}
                \texttt{Penalizace: (7,6) není prázdný, ale obsahuje '─'.} \\
\\ 
\texttt{Dohromady jste nasbírali 14 z 16 pokut.} \\
\texttt{Zkuste to prosím znovu!} \\
            }
        }
    }
    & & \\ \\

    \theutterance \stepcounter{utterance}  
    & & \multicolumn{4}{p{0.6\linewidth}}{
        \cellcolor[rgb]{0.9,0.9,0.9}{
            \makecell[{{p{\linewidth}}}]{
                \texttt{\tiny{[P1$\langle$GM]}}
                \texttt{Penalizace: (7,6) není prázdný, ale obsahuje '─'.} \\
\\ 
\texttt{Dohromady jste nasbírali 14 z 16 pokut.} \\
\texttt{Zkuste to prosím znovu!} \\
            }
        }
    }
    & \\ \\

    \theutterance \stepcounter{utterance}  
    & \multicolumn{4}{p{0.6\linewidth}}{
        \cellcolor[rgb]{0.9,0.9,0.9}{
            \makecell[{{p{\linewidth}}}]{
                \texttt{\tiny{[P1$\rangle$GM]}}
                \texttt{POHYB: D, (4,6)} \\
            }
        }
    }
    & & \\ \\

    \theutterance \stepcounter{utterance}  
    & & & \multicolumn{2}{p{0.3\linewidth}}{
        \cellcolor[rgb]{0.9,0.9,0.9}{
            \makecell[{{p{\linewidth}}}]{
                \texttt{\tiny{[GM$|$GM]}}
                \texttt{POHYB: D, (4,6)} \\
            }
        }
    }
    & & \\ \\

    \theutterance \stepcounter{utterance}  
    & & & \multicolumn{2}{p{0.3\linewidth}}{
        \cellcolor[rgb]{0.9,0.9,0.9}{
            \makecell[{{p{\linewidth}}}]{
                \texttt{\tiny{[GM$|$GM]}}
                \texttt{Penalizace: (4,6) není prázdný, ale obsahuje '┴'.} \\
\\ 
\texttt{Dohromady jste nasbírali 15 z 16 pokut.} \\
\texttt{Zkuste to prosím znovu!} \\
            }
        }
    }
    & & \\ \\

    \theutterance \stepcounter{utterance}  
    & & \multicolumn{4}{p{0.6\linewidth}}{
        \cellcolor[rgb]{0.9,0.9,0.9}{
            \makecell[{{p{\linewidth}}}]{
                \texttt{\tiny{[P1$\langle$GM]}}
                \texttt{Penalizace: (4,6) není prázdný, ale obsahuje '┴'.} \\
\\ 
\texttt{Dohromady jste nasbírali 15 z 16 pokut.} \\
\texttt{Zkuste to prosím znovu!} \\
            }
        }
    }
    & \\ \\

    \theutterance \stepcounter{utterance}  
    & \multicolumn{4}{p{0.6\linewidth}}{
        \cellcolor[rgb]{0.9,0.9,0.9}{
            \makecell[{{p{\linewidth}}}]{
                \texttt{\tiny{[P1$\rangle$GM]}}
                \texttt{POHYB: M, (5,7)} \\
            }
        }
    }
    & & \\ \\

    \theutterance \stepcounter{utterance}  
    & & & \multicolumn{2}{p{0.3\linewidth}}{
        \cellcolor[rgb]{0.9,0.9,0.9}{
            \makecell[{{p{\linewidth}}}]{
                \texttt{\tiny{[GM$|$GM]}}
                \texttt{POHYB: M, (5,7)} \\
            }
        }
    }
    & & \\ \\

    \theutterance \stepcounter{utterance}  
    & & & \multicolumn{2}{p{0.3\linewidth}}{
        \cellcolor[rgb]{0.9,0.9,0.9}{
            \makecell[{{p{\linewidth}}}]{
                \texttt{\tiny{[GM$|$GM]}}
                \texttt{Penalizace: (5,7) není prázdný, ale obsahuje 'M'.} \\
\\ 
\texttt{Dohromady jste nasbírali 16 z 16 pokut.} \\
\texttt{Zkuste to prosím znovu!} \\
            }
        }
    }
    & & \\ \\

    \theutterance \stepcounter{utterance}  
    & & \multicolumn{4}{p{0.6\linewidth}}{
        \cellcolor[rgb]{0.9,0.9,0.9}{
            \makecell[{{p{\linewidth}}}]{
                \texttt{\tiny{[P1$\langle$GM]}}
                \texttt{Penalizace: (5,7) není prázdný, ale obsahuje 'M'.} \\
\\ 
\texttt{Dohromady jste nasbírali 16 z 16 pokut.} \\
\texttt{Zkuste to prosím znovu!} \\
            }
        }
    }
    & \\ \\

    \theutterance \stepcounter{utterance}  
    & \multicolumn{4}{p{0.6\linewidth}}{
        \cellcolor[rgb]{0.9,0.9,0.9}{
            \makecell[{{p{\linewidth}}}]{
                \texttt{\tiny{[P1$\rangle$GM]}}
                \texttt{ŘEKNI: Mám potíže s přesuny, nabral jsem hodně penalizací. Moje aktuální pozice: P (1,1), O (1,3), T (5,2), S (5,5), D (4,7), A (7,1), M (5,7). Potřebuji přesunout D na (6,6) a M na (5,7) už tam je. Jak jsi na tom ty?} \\
            }
        }
    }
    & & \\ \\

    \theutterance \stepcounter{utterance}  
    & & & \multicolumn{2}{p{0.3\linewidth}}{
        \cellcolor[rgb]{0.9,0.9,0.9}{
            \makecell[{{p{\linewidth}}}]{
                \texttt{\tiny{[GM$|$GM]}}
                \texttt{ŘEKNI: Mám potíže s přesuny, nabral jsem hodně penalizací. Moje aktuální pozice: P (1,1), O (1,3), T (5,2), S (5,5), D (4,7), A (7,1), M (5,7). Potřebuji přesunout D na (6,6) a M na (5,7) už tam je. Jak jsi na tom ty?} \\
            }
        }
    }
    & & \\ \\

    \theutterance \stepcounter{utterance}  
    & & \multicolumn{4}{p{0.6\linewidth}}{
        \cellcolor[rgb]{0.9,0.9,0.9}{
            \makecell[{{p{\linewidth}}}]{
                \texttt{\tiny{[GM$\rangle$P2]}}
                \texttt{Úspěšně přesunuto 'M' do (5,7). Vaše aktualizovaná mřížka vypadá takto:} \\
\texttt{```} \\
\texttt{ 1234567} \\
\texttt{╔═╤═════╗} \\
\texttt{║P│◌◌◌◌A║ 1} \\
\texttt{╟─┤T◌◌◌◌║ 2} \\
\texttt{║O└─┐◌◌◌║ 3} \\
\texttt{║◌◌◌├─┐◌║ 4} \\
\texttt{║◌◌◌│S│◌║ 5} \\
\texttt{╟─┬─┴─┼─╢ 6} \\
\texttt{║◌│◌◌M│D║ 7} \\
\texttt{╚═╧═══╧═╝} \\
\texttt{```} \\
\texttt{Právě hrajete kolo 9 z maximálně 28.} \\
\texttt{Dohromady jste nasbírali 16 z 16 pokut.} \\
\texttt{Druhý hráč poslal tuto zprávu:} \\
\\ 
\\ 
\texttt{"Mám potíže s přesuny, nabral jsem hodně penalizací. Moje aktuální pozice: P (1,1), O (1,3), T (5,2), S (5,5), D (4,7), A (7,1), M (5,7). Potřebuji přesunout D na (6,6) a M na (5,7) už tam je. Jak jsi na tom ty?"} \\
\\ 
\\ 
\texttt{Jaký je váš další příkaz?} \\
            }
        }
    }
    & \\ \\

    \theutterance \stepcounter{utterance}  
    & & & \multicolumn{4}{p{0.6\linewidth}}{
        \cellcolor[rgb]{0.9,0.9,0.9}{
            \makecell[{{p{\linewidth}}}]{
                \texttt{\tiny{[GM$\langle$P2]}}
                \texttt{ŘEKNI: Já mám hotovo: P (1,1), O (1,3), T (3,2), S (5,5), M (5,7), A (7,1), D (7,7). Potřebuju přesunout T z (3,2) na (3,3) a D z (7,7) na (6,6). Oba potřebujeme D na (6,6) a T na (3,3). Zbývá nám 0 pokut, tak opatrně!} \\
            }
        }
    }
     \\ \\

    \theutterance \stepcounter{utterance}  
    & & & \multicolumn{2}{p{0.3\linewidth}}{
        \cellcolor[rgb]{0.9,0.9,0.9}{
            \makecell[{{p{\linewidth}}}]{
                \texttt{\tiny{[GM$|$GM]}}
                \texttt{ŘEKNI: Já mám hotovo: P (1,1), O (1,3), T (3,2), S (5,5), M (5,7), A (7,1), D (7,7). Potřebuju přesunout T z (3,2) na (3,3) a D z (7,7) na (6,6). Oba potřebujeme D na (6,6) a T na (3,3). Zbývá nám 0 pokut, tak opatrně!} \\
            }
        }
    }
    & & \\ \\

    \theutterance \stepcounter{utterance}  
    & & \multicolumn{4}{p{0.6\linewidth}}{
        \cellcolor[rgb]{0.9,0.9,0.9}{
            \makecell[{{p{\linewidth}}}]{
                \texttt{\tiny{[P1$\langle$GM]}}
                \texttt{Vaše zpráva byla předána druhému hráči.} \\
\\ 
\\ 
\texttt{Právě hrajete kolo 9 z maximálně 28.} \\
\texttt{Dohromady jste nasbírali 16 z 16 pokut.} \\
\texttt{Druhý hráč poslal tuto zprávu:} \\
\\ 
\\ 
\texttt{"Já mám hotovo: P (1,1), O (1,3), T (3,2), S (5,5), M (5,7), A (7,1), D (7,7). Potřebuju přesunout T z (3,2) na (3,3) a D z (7,7) na (6,6). Oba potřebujeme D na (6,6) a T na (3,3). Zbývá nám 0 pokut, tak opatrně!"} \\
\\ 
\\ 
\texttt{Jaký je váš další příkaz?} \\
            }
        }
    }
    & \\ \\

    \theutterance \stepcounter{utterance}  
    & \multicolumn{4}{p{0.6\linewidth}}{
        \cellcolor[rgb]{0.9,0.9,0.9}{
            \makecell[{{p{\linewidth}}}]{
                \texttt{\tiny{[P1$\rangle$GM]}}
                \texttt{POHYB: D, (3,7)} \\
            }
        }
    }
    & & \\ \\

    \theutterance \stepcounter{utterance}  
    & & & \multicolumn{2}{p{0.3\linewidth}}{
        \cellcolor[rgb]{0.9,0.9,0.9}{
            \makecell[{{p{\linewidth}}}]{
                \texttt{\tiny{[GM$|$GM]}}
                \texttt{POHYB: D, (3,7)} \\
            }
        }
    }
    & & \\ \\

    \theutterance \stepcounter{utterance}  
    & & & \multicolumn{2}{p{0.3\linewidth}}{
        \cellcolor[rgb]{0.9,0.9,0.9}{
            \makecell[{{p{\linewidth}}}]{
                \texttt{\tiny{[GM$|$GM]}}
                \texttt{D moved to (3, 7)} \\
            }
        }
    }
    & & \\ \\

    \theutterance \stepcounter{utterance}  
    & & \multicolumn{4}{p{0.6\linewidth}}{
        \cellcolor[rgb]{0.9,0.9,0.9}{
            \makecell[{{p{\linewidth}}}]{
                \texttt{\tiny{[GM$\rangle$P2]}}
                \texttt{Vaše zpráva byla předána druhému hráči.} \\
\\ 
\\ 
\texttt{Právě hrajete kolo 10 z maximálně 28.} \\
\texttt{Dohromady jste nasbírali 16 z 16 pokut.} \\
\texttt{Druhý hráč přesunul objekt na své mřížce.} \\
\\ 
\\ 
\texttt{Jaký je váš další příkaz?} \\
            }
        }
    }
    & \\ \\

    \theutterance \stepcounter{utterance}  
    & & & \multicolumn{4}{p{0.6\linewidth}}{
        \cellcolor[rgb]{0.9,0.9,0.9}{
            \makecell[{{p{\linewidth}}}]{
                \texttt{\tiny{[GM$\langle$P2]}}
                \texttt{POHYB: T, (3,3)} \\
            }
        }
    }
     \\ \\

    \theutterance \stepcounter{utterance}  
    & & & \multicolumn{2}{p{0.3\linewidth}}{
        \cellcolor[rgb]{0.9,0.9,0.9}{
            \makecell[{{p{\linewidth}}}]{
                \texttt{\tiny{[GM$|$GM]}}
                \texttt{POHYB: T, (3,3)} \\
            }
        }
    }
    & & \\ \\

    \theutterance \stepcounter{utterance}  
    & & & \multicolumn{2}{p{0.3\linewidth}}{
        \cellcolor[rgb]{0.9,0.9,0.9}{
            \makecell[{{p{\linewidth}}}]{
                \texttt{\tiny{[GM$|$GM]}}
                \texttt{Penalizace: (3,3) není prázdný, ale obsahuje '─'.} \\
\\ 
\texttt{Dohromady jste nasbírali 17 z 16 pokut.} \\
\texttt{Zkuste to prosím znovu!} \\
            }
        }
    }
    & & \\ \\

    \theutterance \stepcounter{utterance}  
    & & \multicolumn{4}{p{0.6\linewidth}}{
        \cellcolor[rgb]{0.9,0.9,0.9}{
            \makecell[{{p{\linewidth}}}]{
                \texttt{\tiny{[GM$\rangle$P2]}}
                \texttt{Penalizace: (3,3) není prázdný, ale obsahuje '─'.} \\
\\ 
\texttt{Dohromady jste nasbírali 17 z 16 pokut.} \\
\texttt{Zkuste to prosím znovu!} \\
            }
        }
    }
    & \\ \\

    \theutterance \stepcounter{utterance}  
    & & & \multicolumn{4}{p{0.6\linewidth}}{
        \cellcolor[rgb]{0.9,0.9,0.9}{
            \makecell[{{p{\linewidth}}}]{
                \texttt{\tiny{[GM$\langle$P2]}}
                \texttt{POHYB: T, (3,4)} \\
            }
        }
    }
     \\ \\

    \theutterance \stepcounter{utterance}  
    & & & \multicolumn{2}{p{0.3\linewidth}}{
        \cellcolor[rgb]{0.9,0.9,0.9}{
            \makecell[{{p{\linewidth}}}]{
                \texttt{\tiny{[GM$|$GM]}}
                \texttt{POHYB: T, (3,4)} \\
            }
        }
    }
    & & \\ \\

    \theutterance \stepcounter{utterance}  
    & & & \multicolumn{2}{p{0.3\linewidth}}{
        \cellcolor[rgb]{0.9,0.9,0.9}{
            \makecell[{{p{\linewidth}}}]{
                \texttt{\tiny{[GM$|$GM]}}
                \texttt{T moved to (3, 4)} \\
            }
        }
    }
    & & \\ \\

    \theutterance \stepcounter{utterance}  
    & & \multicolumn{4}{p{0.6\linewidth}}{
        \cellcolor[rgb]{0.9,0.9,0.9}{
            \makecell[{{p{\linewidth}}}]{
                \texttt{\tiny{[P1$\langle$GM]}}
                \texttt{Úspěšně přesunuto 'D' do (3,7). Vaše aktualizovaná mřížka vypadá takto:} \\
\texttt{```} \\
\texttt{ 1234567} \\
\texttt{╔═╤═════╗} \\
\texttt{║P│◌◌◌◌A║ 1} \\
\texttt{╟─┤◌◌T◌◌║ 2} \\
\texttt{║O└─┐◌◌◌║ 3} \\
\texttt{║◌◌◌├─┐◌║ 4} \\
\texttt{║◌◌◌│S│◌║ 5} \\
\texttt{╟─┬─┴─┼─╢ 6} \\
\texttt{║◌│D◌M│◌║ 7} \\
\texttt{╚═╧═══╧═╝} \\
\texttt{```} \\
\texttt{Právě hrajete kolo 10 z maximálně 28.} \\
\texttt{Dohromady jste nasbírali 17 z 16 pokut.} \\
\texttt{Druhý hráč přesunul objekt na své mřížce.} \\
\\ 
\\ 
\texttt{Jaký je váš další příkaz?} \\
            }
        }
    }
    & \\ \\

    \theutterance \stepcounter{utterance}  
    & \multicolumn{4}{p{0.6\linewidth}}{
        \cellcolor[rgb]{0.9,0.9,0.9}{
            \makecell[{{p{\linewidth}}}]{
                \texttt{\tiny{[P1$\rangle$GM]}}
                \texttt{POHYB: D, (3,6)} \\
            }
        }
    }
    & & \\ \\

    \theutterance \stepcounter{utterance}  
    & & & \multicolumn{2}{p{0.3\linewidth}}{
        \cellcolor[rgb]{0.9,0.9,0.9}{
            \makecell[{{p{\linewidth}}}]{
                \texttt{\tiny{[GM$|$GM]}}
                \texttt{POHYB: D, (3,6)} \\
            }
        }
    }
    & & \\ \\

    \theutterance \stepcounter{utterance}  
    & & & \multicolumn{2}{p{0.3\linewidth}}{
        \cellcolor[rgb]{0.9,0.9,0.9}{
            \makecell[{{p{\linewidth}}}]{
                \texttt{\tiny{[GM$|$GM]}}
                \texttt{Penalizace: (3,6) není prázdný, ale obsahuje '─'.} \\
\\ 
\texttt{Dohromady jste nasbírali 18 z 16 pokut.} \\
\texttt{Zkuste to prosím znovu!} \\
            }
        }
    }
    & & \\ \\

    \theutterance \stepcounter{utterance}  
    & & \multicolumn{4}{p{0.6\linewidth}}{
        \cellcolor[rgb]{0.9,0.9,0.9}{
            \makecell[{{p{\linewidth}}}]{
                \texttt{\tiny{[P1$\langle$GM]}}
                \texttt{Penalizace: (3,6) není prázdný, ale obsahuje '─'.} \\
\\ 
\texttt{Dohromady jste nasbírali 18 z 16 pokut.} \\
\texttt{Zkuste to prosím znovu!} \\
            }
        }
    }
    & \\ \\

    \theutterance \stepcounter{utterance}  
    & \multicolumn{4}{p{0.6\linewidth}}{
        \cellcolor[rgb]{0.9,0.9,0.9}{
            \makecell[{{p{\linewidth}}}]{
                \texttt{\tiny{[P1$\rangle$GM]}}
                \texttt{POHYB: D, (2,6)} \\
            }
        }
    }
    & & \\ \\

    \theutterance \stepcounter{utterance}  
    & & & \multicolumn{2}{p{0.3\linewidth}}{
        \cellcolor[rgb]{0.9,0.9,0.9}{
            \makecell[{{p{\linewidth}}}]{
                \texttt{\tiny{[GM$|$GM]}}
                \texttt{POHYB: D, (2,6)} \\
            }
        }
    }
    & & \\ \\

    \theutterance \stepcounter{utterance}  
    & & & \multicolumn{2}{p{0.3\linewidth}}{
        \cellcolor[rgb]{0.9,0.9,0.9}{
            \makecell[{{p{\linewidth}}}]{
                \texttt{\tiny{[GM$|$GM]}}
                \texttt{Penalizace: (2,6) není prázdný, ale obsahuje '┬'.} \\
\\ 
\texttt{Dohromady jste nasbírali 19 z 16 pokut.} \\
\texttt{Zkuste to prosím znovu!} \\
            }
        }
    }
    & & \\ \\

    \theutterance \stepcounter{utterance}  
    & & \multicolumn{4}{p{0.6\linewidth}}{
        \cellcolor[rgb]{0.9,0.9,0.9}{
            \makecell[{{p{\linewidth}}}]{
                \texttt{\tiny{[P1$\langle$GM]}}
                \texttt{Penalizace: (2,6) není prázdný, ale obsahuje '┬'.} \\
\\ 
\texttt{Dohromady jste nasbírali 19 z 16 pokut.} \\
\texttt{Zkuste to prosím znovu!} \\
            }
        }
    }
    & \\ \\

    \theutterance \stepcounter{utterance}  
    & \multicolumn{4}{p{0.6\linewidth}}{
        \cellcolor[rgb]{0.9,0.9,0.9}{
            \makecell[{{p{\linewidth}}}]{
                \texttt{\tiny{[P1$\rangle$GM]}}
                \texttt{POHYB: D, (1,6)} \\
            }
        }
    }
    & & \\ \\

    \theutterance \stepcounter{utterance}  
    & & & \multicolumn{2}{p{0.3\linewidth}}{
        \cellcolor[rgb]{0.9,0.9,0.9}{
            \makecell[{{p{\linewidth}}}]{
                \texttt{\tiny{[GM$|$GM]}}
                \texttt{POHYB: D, (1,6)} \\
            }
        }
    }
    & & \\ \\

    \theutterance \stepcounter{utterance}  
    & & & \multicolumn{2}{p{0.3\linewidth}}{
        \cellcolor[rgb]{0.9,0.9,0.9}{
            \makecell[{{p{\linewidth}}}]{
                \texttt{\tiny{[GM$|$GM]}}
                \texttt{Penalizace: (1,6) není prázdný, ale obsahuje '─'.} \\
\\ 
\texttt{Dohromady jste nasbírali 20 z 16 pokut.} \\
\texttt{Zkuste to prosím znovu!} \\
            }
        }
    }
    & & \\ \\

    \theutterance \stepcounter{utterance}  
    & & \multicolumn{4}{p{0.6\linewidth}}{
        \cellcolor[rgb]{0.9,0.9,0.9}{
            \makecell[{{p{\linewidth}}}]{
                \texttt{\tiny{[P1$\langle$GM]}}
                \texttt{Penalizace: (1,6) není prázdný, ale obsahuje '─'.} \\
\\ 
\texttt{Dohromady jste nasbírali 20 z 16 pokut.} \\
\texttt{Zkuste to prosím znovu!} \\
            }
        }
    }
    & \\ \\

    \theutterance \stepcounter{utterance}  
    & \multicolumn{4}{p{0.6\linewidth}}{
        \cellcolor[rgb]{0.9,0.9,0.9}{
            \makecell[{{p{\linewidth}}}]{
                \texttt{\tiny{[P1$\rangle$GM]}}
                \texttt{POHYB: D, (3,4)} \\
            }
        }
    }
    & & \\ \\

    \theutterance \stepcounter{utterance}  
    & & & \multicolumn{2}{p{0.3\linewidth}}{
        \cellcolor[rgb]{0.9,0.9,0.9}{
            \makecell[{{p{\linewidth}}}]{
                \texttt{\tiny{[GM$|$GM]}}
                \texttt{POHYB: D, (3,4)} \\
            }
        }
    }
    & & \\ \\

    \theutterance \stepcounter{utterance}  
    & & & \multicolumn{2}{p{0.3\linewidth}}{
        \cellcolor[rgb]{0.9,0.9,0.9}{
            \makecell[{{p{\linewidth}}}]{
                \texttt{\tiny{[GM$|$GM]}}
                \texttt{D moved to (3, 4)} \\
            }
        }
    }
    & & \\ \\

    \theutterance \stepcounter{utterance}  
    & & \multicolumn{4}{p{0.6\linewidth}}{
        \cellcolor[rgb]{0.9,0.9,0.9}{
            \makecell[{{p{\linewidth}}}]{
                \texttt{\tiny{[GM$\rangle$P2]}}
                \texttt{Úspěšně přesunuto 'T' do (3,4). Vaše aktualizovaná mřížka vypadá takto:} \\
\texttt{```} \\
\texttt{ 1234567} \\
\texttt{╔═╤═════╗} \\
\texttt{║P│◌◌◌◌A║ 1} \\
\texttt{╟─┤◌◌◌◌◌║ 2} \\
\texttt{║O└─┐◌◌◌║ 3} \\
\texttt{║◌◌T├─┐◌║ 4} \\
\texttt{║◌◌◌│S│◌║ 5} \\
\texttt{╟─┬─┴─┼─╢ 6} \\
\texttt{║◌│◌◌M│D║ 7} \\
\texttt{╚═╧═══╧═╝} \\
\texttt{```} \\
\texttt{Právě hrajete kolo 11 z maximálně 28.} \\
\texttt{Dohromady jste nasbírali 20 z 16 pokut.} \\
\texttt{Druhý hráč přesunul objekt na své mřížce.} \\
\\ 
\\ 
\texttt{Jaký je váš další příkaz?} \\
            }
        }
    }
    & \\ \\

    \theutterance \stepcounter{utterance}  
    & & & \multicolumn{4}{p{0.6\linewidth}}{
        \cellcolor[rgb]{0.9,0.9,0.9}{
            \makecell[{{p{\linewidth}}}]{
                \texttt{\tiny{[GM$\langle$P2]}}
                \texttt{POHYB: D, (6,6)} \\
            }
        }
    }
     \\ \\

    \theutterance \stepcounter{utterance}  
    & & & \multicolumn{2}{p{0.3\linewidth}}{
        \cellcolor[rgb]{0.9,0.9,0.9}{
            \makecell[{{p{\linewidth}}}]{
                \texttt{\tiny{[GM$|$GM]}}
                \texttt{POHYB: D, (6,6)} \\
            }
        }
    }
    & & \\ \\

    \theutterance \stepcounter{utterance}  
    & & & \multicolumn{2}{p{0.3\linewidth}}{
        \cellcolor[rgb]{0.9,0.9,0.9}{
            \makecell[{{p{\linewidth}}}]{
                \texttt{\tiny{[GM$|$GM]}}
                \texttt{Penalizace: (6,6) není prázdný, ale obsahuje '┼'.} \\
\\ 
\texttt{Dohromady jste nasbírali 21 z 16 pokut.} \\
\texttt{Zkuste to prosím znovu!} \\
            }
        }
    }
    & & \\ \\

    \theutterance \stepcounter{utterance}  
    & & \multicolumn{4}{p{0.6\linewidth}}{
        \cellcolor[rgb]{0.9,0.9,0.9}{
            \makecell[{{p{\linewidth}}}]{
                \texttt{\tiny{[GM$\rangle$P2]}}
                \texttt{Penalizace: (6,6) není prázdný, ale obsahuje '┼'.} \\
\\ 
\texttt{Dohromady jste nasbírali 21 z 16 pokut.} \\
\texttt{Zkuste to prosím znovu!} \\
            }
        }
    }
    & \\ \\

    \theutterance \stepcounter{utterance}  
    & & & \multicolumn{4}{p{0.6\linewidth}}{
        \cellcolor[rgb]{0.9,0.9,0.9}{
            \makecell[{{p{\linewidth}}}]{
                \texttt{\tiny{[GM$\langle$P2]}}
                \texttt{POHYB: D, (5,6)} \\
            }
        }
    }
     \\ \\

    \theutterance \stepcounter{utterance}  
    & & & \multicolumn{2}{p{0.3\linewidth}}{
        \cellcolor[rgb]{0.9,0.9,0.9}{
            \makecell[{{p{\linewidth}}}]{
                \texttt{\tiny{[GM$|$GM]}}
                \texttt{POHYB: D, (5,6)} \\
            }
        }
    }
    & & \\ \\

    \theutterance \stepcounter{utterance}  
    & & & \multicolumn{2}{p{0.3\linewidth}}{
        \cellcolor[rgb]{0.9,0.9,0.9}{
            \makecell[{{p{\linewidth}}}]{
                \texttt{\tiny{[GM$|$GM]}}
                \texttt{Penalizace: (5,6) není prázdný, ale obsahuje '─'.} \\
\\ 
\texttt{Dohromady jste nasbírali 22 z 16 pokut.} \\
\texttt{Zkuste to prosím znovu!} \\
            }
        }
    }
    & & \\ \\

    \theutterance \stepcounter{utterance}  
    & & \multicolumn{4}{p{0.6\linewidth}}{
        \cellcolor[rgb]{0.9,0.9,0.9}{
            \makecell[{{p{\linewidth}}}]{
                \texttt{\tiny{[GM$\rangle$P2]}}
                \texttt{Penalizace: (5,6) není prázdný, ale obsahuje '─'.} \\
\\ 
\texttt{Dohromady jste nasbírali 22 z 16 pokut.} \\
\texttt{Zkuste to prosím znovu!} \\
            }
        }
    }
    & \\ \\

    \theutterance \stepcounter{utterance}  
    & & & \multicolumn{4}{p{0.6\linewidth}}{
        \cellcolor[rgb]{0.9,0.9,0.9}{
            \makecell[{{p{\linewidth}}}]{
                \texttt{\tiny{[GM$\langle$P2]}}
                \texttt{POHYB: D, (4,6)} \\
            }
        }
    }
     \\ \\

    \theutterance \stepcounter{utterance}  
    & & & \multicolumn{2}{p{0.3\linewidth}}{
        \cellcolor[rgb]{0.9,0.9,0.9}{
            \makecell[{{p{\linewidth}}}]{
                \texttt{\tiny{[GM$|$GM]}}
                \texttt{POHYB: D, (4,6)} \\
            }
        }
    }
    & & \\ \\

    \theutterance \stepcounter{utterance}  
    & & & \multicolumn{2}{p{0.3\linewidth}}{
        \cellcolor[rgb]{0.9,0.9,0.9}{
            \makecell[{{p{\linewidth}}}]{
                \texttt{\tiny{[GM$|$GM]}}
                \texttt{Penalizace: (4,6) není prázdný, ale obsahuje '┴'.} \\
\\ 
\texttt{Dohromady jste nasbírali 23 z 16 pokut.} \\
\texttt{Zkuste to prosím znovu!} \\
            }
        }
    }
    & & \\ \\

    \theutterance \stepcounter{utterance}  
    & & \multicolumn{4}{p{0.6\linewidth}}{
        \cellcolor[rgb]{0.9,0.9,0.9}{
            \makecell[{{p{\linewidth}}}]{
                \texttt{\tiny{[GM$\rangle$P2]}}
                \texttt{Penalizace: (4,6) není prázdný, ale obsahuje '┴'.} \\
\\ 
\texttt{Dohromady jste nasbírali 23 z 16 pokut.} \\
\texttt{Zkuste to prosím znovu!} \\
            }
        }
    }
    & \\ \\

    \theutterance \stepcounter{utterance}  
    & & & \multicolumn{4}{p{0.6\linewidth}}{
        \cellcolor[rgb]{0.9,0.9,0.9}{
            \makecell[{{p{\linewidth}}}]{
                \texttt{\tiny{[GM$\langle$P2]}}
                \texttt{POHYB: D, (6,5)} \\
            }
        }
    }
     \\ \\

    \theutterance \stepcounter{utterance}  
    & & & \multicolumn{2}{p{0.3\linewidth}}{
        \cellcolor[rgb]{0.9,0.9,0.9}{
            \makecell[{{p{\linewidth}}}]{
                \texttt{\tiny{[GM$|$GM]}}
                \texttt{POHYB: D, (6,5)} \\
            }
        }
    }
    & & \\ \\

    \theutterance \stepcounter{utterance}  
    & & & \multicolumn{2}{p{0.3\linewidth}}{
        \cellcolor[rgb]{0.9,0.9,0.9}{
            \makecell[{{p{\linewidth}}}]{
                \texttt{\tiny{[GM$|$GM]}}
                \texttt{Penalizace: (6,5) není prázdný, ale obsahuje '│'.} \\
\\ 
\texttt{Dohromady jste nasbírali 24 z 16 pokut.} \\
\texttt{Zkuste to prosím znovu!} \\
            }
        }
    }
    & & \\ \\

    \theutterance \stepcounter{utterance}  
    & & \multicolumn{4}{p{0.6\linewidth}}{
        \cellcolor[rgb]{0.9,0.9,0.9}{
            \makecell[{{p{\linewidth}}}]{
                \texttt{\tiny{[GM$\rangle$P2]}}
                \texttt{Penalizace: (6,5) není prázdný, ale obsahuje '│'.} \\
\\ 
\texttt{Dohromady jste nasbírali 24 z 16 pokut.} \\
\texttt{Zkuste to prosím znovu!} \\
            }
        }
    }
    & \\ \\

    \theutterance \stepcounter{utterance}  
    & & & \multicolumn{4}{p{0.6\linewidth}}{
        \cellcolor[rgb]{0.9,0.9,0.9}{
            \makecell[{{p{\linewidth}}}]{
                \texttt{\tiny{[GM$\langle$P2]}}
                \texttt{ŘEKNI: Nemůžu přesunout D na (6,6), je tam překážka. Moje mřížka má na (6,6) symbol '┼'. Můžeme změnit cíl pro D? Co třeba D na (7,7) nebo jinam? Kde máš volno?} \\
            }
        }
    }
     \\ \\

    \theutterance \stepcounter{utterance}  
    & & & \multicolumn{2}{p{0.3\linewidth}}{
        \cellcolor[rgb]{0.9,0.9,0.9}{
            \makecell[{{p{\linewidth}}}]{
                \texttt{\tiny{[GM$|$GM]}}
                \texttt{ŘEKNI: Nemůžu přesunout D na (6,6), je tam překážka. Moje mřížka má na (6,6) symbol '┼'. Můžeme změnit cíl pro D? Co třeba D na (7,7) nebo jinam? Kde máš volno?} \\
            }
        }
    }
    & & \\ \\

    \theutterance \stepcounter{utterance}  
    & & \multicolumn{4}{p{0.6\linewidth}}{
        \cellcolor[rgb]{0.9,0.9,0.9}{
            \makecell[{{p{\linewidth}}}]{
                \texttt{\tiny{[P1$\langle$GM]}}
                \texttt{Úspěšně přesunuto 'D' do (3,4). Vaše aktualizovaná mřížka vypadá takto:} \\
\texttt{```} \\
\texttt{ 1234567} \\
\texttt{╔═╤═════╗} \\
\texttt{║P│◌◌◌◌A║ 1} \\
\texttt{╟─┤◌◌T◌◌║ 2} \\
\texttt{║O└─┐◌◌◌║ 3} \\
\texttt{║◌◌D├─┐◌║ 4} \\
\texttt{║◌◌◌│S│◌║ 5} \\
\texttt{╟─┬─┴─┼─╢ 6} \\
\texttt{║◌│◌◌M│◌║ 7} \\
\texttt{╚═╧═══╧═╝} \\
\texttt{```} \\
\texttt{Právě hrajete kolo 11 z maximálně 28.} \\
\texttt{Dohromady jste nasbírali 24 z 16 pokut.} \\
\texttt{Druhý hráč poslal tuto zprávu:} \\
\\ 
\\ 
\texttt{"Nemůžu přesunout D na (6,6), je tam překážka. Moje mřížka má na (6,6) symbol '┼'. Můžeme změnit cíl pro D? Co třeba D na (7,7) nebo jinam? Kde máš volno?"} \\
\\ 
\\ 
\texttt{Jaký je váš další příkaz?} \\
            }
        }
    }
    & \\ \\

    \theutterance \stepcounter{utterance}  
    & \multicolumn{4}{p{0.6\linewidth}}{
        \cellcolor[rgb]{0.9,0.9,0.9}{
            \makecell[{{p{\linewidth}}}]{
                \texttt{\tiny{[P1$\rangle$GM]}}
                \texttt{ŘEKNI: Souhlasím se změnou cíle pro D. U mě je volné (7,7). Můžeme dát D na (7,7). Také T ještě nemám na (3,3), mám ho na (5,2). Jaké máš volné pozice kolem (3,3)?} \\
            }
        }
    }
    & & \\ \\

    \theutterance \stepcounter{utterance}  
    & & & \multicolumn{2}{p{0.3\linewidth}}{
        \cellcolor[rgb]{0.9,0.9,0.9}{
            \makecell[{{p{\linewidth}}}]{
                \texttt{\tiny{[GM$|$GM]}}
                \texttt{ŘEKNI: Souhlasím se změnou cíle pro D. U mě je volné (7,7). Můžeme dát D na (7,7). Také T ještě nemám na (3,3), mám ho na (5,2). Jaké máš volné pozice kolem (3,3)?} \\
            }
        }
    }
    & & \\ \\

    \theutterance \stepcounter{utterance}  
    & & \multicolumn{4}{p{0.6\linewidth}}{
        \cellcolor[rgb]{0.9,0.9,0.9}{
            \makecell[{{p{\linewidth}}}]{
                \texttt{\tiny{[GM$\rangle$P2]}}
                \texttt{Vaše zpráva byla předána druhému hráči.} \\
\\ 
\\ 
\texttt{Právě hrajete kolo 12 z maximálně 28.} \\
\texttt{Dohromady jste nasbírali 24 z 16 pokut.} \\
\texttt{Druhý hráč poslal tuto zprávu:} \\
\\ 
\\ 
\texttt{"Souhlasím se změnou cíle pro D. U mě je volné (7,7). Můžeme dát D na (7,7). Také T ještě nemám na (3,3), mám ho na (5,2). Jaké máš volné pozice kolem (3,3)?"} \\
\\ 
\\ 
\texttt{Jaký je váš další příkaz?} \\
            }
        }
    }
    & \\ \\

    \theutterance \stepcounter{utterance}  
    & & & \multicolumn{4}{p{0.6\linewidth}}{
        \cellcolor[rgb]{0.9,0.9,0.9}{
            \makecell[{{p{\linewidth}}}]{
                \texttt{\tiny{[GM$\langle$P2]}}
                \texttt{ŘEKNI: D na (7,7) mi vyhovuje, už ho tam mám. Pro T: u mě je (3,3) obsazené '─', ale mám T na (3,4). Můžeme cílit T na (3,4)? Nebo zkusit (2,3) či (4,3)?} \\
            }
        }
    }
     \\ \\

    \theutterance \stepcounter{utterance}  
    & & & \multicolumn{2}{p{0.3\linewidth}}{
        \cellcolor[rgb]{0.9,0.9,0.9}{
            \makecell[{{p{\linewidth}}}]{
                \texttt{\tiny{[GM$|$GM]}}
                \texttt{ŘEKNI: D na (7,7) mi vyhovuje, už ho tam mám. Pro T: u mě je (3,3) obsazené '─', ale mám T na (3,4). Můžeme cílit T na (3,4)? Nebo zkusit (2,3) či (4,3)?} \\
            }
        }
    }
    & & \\ \\

    \theutterance \stepcounter{utterance}  
    & & \multicolumn{4}{p{0.6\linewidth}}{
        \cellcolor[rgb]{0.9,0.9,0.9}{
            \makecell[{{p{\linewidth}}}]{
                \texttt{\tiny{[P1$\langle$GM]}}
                \texttt{Vaše zpráva byla předána druhému hráči.} \\
\\ 
\\ 
\texttt{Právě hrajete kolo 12 z maximálně 28.} \\
\texttt{Dohromady jste nasbírali 24 z 16 pokut.} \\
\texttt{Druhý hráč poslal tuto zprávu:} \\
\\ 
\\ 
\texttt{"D na (7,7) mi vyhovuje, už ho tam mám. Pro T: u mě je (3,3) obsazené '─', ale mám T na (3,4). Můžeme cílit T na (3,4)? Nebo zkusit (2,3) či (4,3)?"} \\
\\ 
\\ 
\texttt{Jaký je váš další příkaz?} \\
            }
        }
    }
    & \\ \\

    \theutterance \stepcounter{utterance}  
    & \multicolumn{4}{p{0.6\linewidth}}{
        \cellcolor[rgb]{0.9,0.9,0.9}{
            \makecell[{{p{\linewidth}}}]{
                \texttt{\tiny{[P1$\rangle$GM]}}
                \texttt{ŘEKNI: U mě je (3,4) obsazené D, takže T na (3,4) teď nejde. Navrhuji přesunout T na (4,2) {-} u mě je to volné. Nebo můžeme zkusit (6,2). Které z těchto máš volné?} \\
            }
        }
    }
    & & \\ \\

    \theutterance \stepcounter{utterance}  
    & & & \multicolumn{2}{p{0.3\linewidth}}{
        \cellcolor[rgb]{0.9,0.9,0.9}{
            \makecell[{{p{\linewidth}}}]{
                \texttt{\tiny{[GM$|$GM]}}
                \texttt{ŘEKNI: U mě je (3,4) obsazené D, takže T na (3,4) teď nejde. Navrhuji přesunout T na (4,2) {-} u mě je to volné. Nebo můžeme zkusit (6,2). Které z těchto máš volné?} \\
            }
        }
    }
    & & \\ \\

    \theutterance \stepcounter{utterance}  
    & & \multicolumn{4}{p{0.6\linewidth}}{
        \cellcolor[rgb]{0.9,0.9,0.9}{
            \makecell[{{p{\linewidth}}}]{
                \texttt{\tiny{[GM$\rangle$P2]}}
                \texttt{Vaše zpráva byla předána druhému hráči.} \\
\\ 
\\ 
\texttt{Právě hrajete kolo 13 z maximálně 28.} \\
\texttt{Dohromady jste nasbírali 24 z 16 pokut.} \\
\texttt{Druhý hráč poslal tuto zprávu:} \\
\\ 
\\ 
\texttt{"U mě je (3,4) obsazené D, takže T na (3,4) teď nejde. Navrhuji přesunout T na (4,2) {-} u mě je to volné. Nebo můžeme zkusit (6,2). Které z těchto máš volné?"} \\
\\ 
\\ 
\texttt{Jaký je váš další příkaz?} \\
            }
        }
    }
    & \\ \\

    \theutterance \stepcounter{utterance}  
    & & & \multicolumn{4}{p{0.6\linewidth}}{
        \cellcolor[rgb]{0.9,0.9,0.9}{
            \makecell[{{p{\linewidth}}}]{
                \texttt{\tiny{[GM$\langle$P2]}}
                \texttt{ŘEKNI: Mám volné (4,2) i (6,2). Souhlasím s T na (4,2). Přesouvám.} \\
            }
        }
    }
     \\ \\

    \theutterance \stepcounter{utterance}  
    & & & \multicolumn{2}{p{0.3\linewidth}}{
        \cellcolor[rgb]{0.9,0.9,0.9}{
            \makecell[{{p{\linewidth}}}]{
                \texttt{\tiny{[GM$|$GM]}}
                \texttt{ŘEKNI: Mám volné (4,2) i (6,2). Souhlasím s T na (4,2). Přesouvám.} \\
            }
        }
    }
    & & \\ \\

    \theutterance \stepcounter{utterance}  
    & & \multicolumn{4}{p{0.6\linewidth}}{
        \cellcolor[rgb]{0.9,0.9,0.9}{
            \makecell[{{p{\linewidth}}}]{
                \texttt{\tiny{[P1$\langle$GM]}}
                \texttt{Vaše zpráva byla předána druhému hráči.} \\
\\ 
\\ 
\texttt{Právě hrajete kolo 13 z maximálně 28.} \\
\texttt{Dohromady jste nasbírali 24 z 16 pokut.} \\
\texttt{Druhý hráč poslal tuto zprávu:} \\
\\ 
\\ 
\texttt{"Mám volné (4,2) i (6,2). Souhlasím s T na (4,2). Přesouvám."} \\
\\ 
\\ 
\texttt{Jaký je váš další příkaz?} \\
            }
        }
    }
    & \\ \\

    \theutterance \stepcounter{utterance}  
    & \multicolumn{4}{p{0.6\linewidth}}{
        \cellcolor[rgb]{0.9,0.9,0.9}{
            \makecell[{{p{\linewidth}}}]{
                \texttt{\tiny{[P1$\rangle$GM]}}
                \texttt{POHYB: T, (4,2)} \\
            }
        }
    }
    & & \\ \\

    \theutterance \stepcounter{utterance}  
    & & & \multicolumn{2}{p{0.3\linewidth}}{
        \cellcolor[rgb]{0.9,0.9,0.9}{
            \makecell[{{p{\linewidth}}}]{
                \texttt{\tiny{[GM$|$GM]}}
                \texttt{POHYB: T, (4,2)} \\
            }
        }
    }
    & & \\ \\

    \theutterance \stepcounter{utterance}  
    & & & \multicolumn{2}{p{0.3\linewidth}}{
        \cellcolor[rgb]{0.9,0.9,0.9}{
            \makecell[{{p{\linewidth}}}]{
                \texttt{\tiny{[GM$|$GM]}}
                \texttt{T moved to (4, 2)} \\
            }
        }
    }
    & & \\ \\

    \theutterance \stepcounter{utterance}  
    & & \multicolumn{4}{p{0.6\linewidth}}{
        \cellcolor[rgb]{0.9,0.9,0.9}{
            \makecell[{{p{\linewidth}}}]{
                \texttt{\tiny{[GM$\rangle$P2]}}
                \texttt{Vaše zpráva byla předána druhému hráči.} \\
\\ 
\\ 
\texttt{Právě hrajete kolo 14 z maximálně 28.} \\
\texttt{Dohromady jste nasbírali 24 z 16 pokut.} \\
\texttt{Druhý hráč přesunul objekt na své mřížce.} \\
\\ 
\\ 
\texttt{Jaký je váš další příkaz?} \\
            }
        }
    }
    & \\ \\

    \theutterance \stepcounter{utterance}  
    & & & \multicolumn{4}{p{0.6\linewidth}}{
        \cellcolor[rgb]{0.9,0.9,0.9}{
            \makecell[{{p{\linewidth}}}]{
                \texttt{\tiny{[GM$\langle$P2]}}
                \texttt{POHYB: T, (4,2)} \\
            }
        }
    }
     \\ \\

    \theutterance \stepcounter{utterance}  
    & & & \multicolumn{2}{p{0.3\linewidth}}{
        \cellcolor[rgb]{0.9,0.9,0.9}{
            \makecell[{{p{\linewidth}}}]{
                \texttt{\tiny{[GM$|$GM]}}
                \texttt{POHYB: T, (4,2)} \\
            }
        }
    }
    & & \\ \\

    \theutterance \stepcounter{utterance}  
    & & & \multicolumn{2}{p{0.3\linewidth}}{
        \cellcolor[rgb]{0.9,0.9,0.9}{
            \makecell[{{p{\linewidth}}}]{
                \texttt{\tiny{[GM$|$GM]}}
                \texttt{T moved to (4, 2)} \\
            }
        }
    }
    & & \\ \\

    \theutterance \stepcounter{utterance}  
    & & \multicolumn{4}{p{0.6\linewidth}}{
        \cellcolor[rgb]{0.9,0.9,0.9}{
            \makecell[{{p{\linewidth}}}]{
                \texttt{\tiny{[P1$\langle$GM]}}
                \texttt{Úspěšně přesunuto 'T' do (4,2). Vaše aktualizovaná mřížka vypadá takto:} \\
\texttt{```} \\
\texttt{ 1234567} \\
\texttt{╔═╤═════╗} \\
\texttt{║P│◌◌◌◌A║ 1} \\
\texttt{╟─┤◌T◌◌◌║ 2} \\
\texttt{║O└─┐◌◌◌║ 3} \\
\texttt{║◌◌D├─┐◌║ 4} \\
\texttt{║◌◌◌│S│◌║ 5} \\
\texttt{╟─┬─┴─┼─╢ 6} \\
\texttt{║◌│◌◌M│◌║ 7} \\
\texttt{╚═╧═══╧═╝} \\
\texttt{```} \\
\texttt{Právě hrajete kolo 14 z maximálně 28.} \\
\texttt{Dohromady jste nasbírali 24 z 16 pokut.} \\
\texttt{Druhý hráč přesunul objekt na své mřížce.} \\
\\ 
\\ 
\texttt{Jaký je váš další příkaz?} \\
            }
        }
    }
    & \\ \\

    \theutterance \stepcounter{utterance}  
    & \multicolumn{4}{p{0.6\linewidth}}{
        \cellcolor[rgb]{0.9,0.9,0.9}{
            \makecell[{{p{\linewidth}}}]{
                \texttt{\tiny{[P1$\rangle$GM]}}
                \texttt{POHYB: D, (7,7)} \\
            }
        }
    }
    & & \\ \\

    \theutterance \stepcounter{utterance}  
    & & & \multicolumn{2}{p{0.3\linewidth}}{
        \cellcolor[rgb]{0.9,0.9,0.9}{
            \makecell[{{p{\linewidth}}}]{
                \texttt{\tiny{[GM$|$GM]}}
                \texttt{POHYB: D, (7,7)} \\
            }
        }
    }
    & & \\ \\

    \theutterance \stepcounter{utterance}  
    & & & \multicolumn{2}{p{0.3\linewidth}}{
        \cellcolor[rgb]{0.9,0.9,0.9}{
            \makecell[{{p{\linewidth}}}]{
                \texttt{\tiny{[GM$|$GM]}}
                \texttt{D moved to (7, 7)} \\
            }
        }
    }
    & & \\ \\

    \theutterance \stepcounter{utterance}  
    & & \multicolumn{4}{p{0.6\linewidth}}{
        \cellcolor[rgb]{0.9,0.9,0.9}{
            \makecell[{{p{\linewidth}}}]{
                \texttt{\tiny{[GM$\rangle$P2]}}
                \texttt{Úspěšně přesunuto 'T' do (4,2). Vaše aktualizovaná mřížka vypadá takto:} \\
\texttt{```} \\
\texttt{ 1234567} \\
\texttt{╔═╤═════╗} \\
\texttt{║P│◌◌◌◌A║ 1} \\
\texttt{╟─┤◌T◌◌◌║ 2} \\
\texttt{║O└─┐◌◌◌║ 3} \\
\texttt{║◌◌◌├─┐◌║ 4} \\
\texttt{║◌◌◌│S│◌║ 5} \\
\texttt{╟─┬─┴─┼─╢ 6} \\
\texttt{║◌│◌◌M│D║ 7} \\
\texttt{╚═╧═══╧═╝} \\
\texttt{```} \\
\texttt{Právě hrajete kolo 15 z maximálně 28.} \\
\texttt{Dohromady jste nasbírali 24 z 16 pokut.} \\
\texttt{Druhý hráč přesunul objekt na své mřížce.} \\
\\ 
\\ 
\texttt{Jaký je váš další příkaz?} \\
            }
        }
    }
    & \\ \\

    \theutterance \stepcounter{utterance}  
    & & & \multicolumn{4}{p{0.6\linewidth}}{
        \cellcolor[rgb]{0.9,0.9,0.9}{
            \makecell[{{p{\linewidth}}}]{
                \texttt{\tiny{[GM$\langle$P2]}}
                \texttt{ŘEKNI: Moje aktuální pozice: P (1,1), O (1,3), T (4,2), S (5,5), M (5,7), A (7,1), D (7,7). Všechny objekty mám na cílových pozicích kromě T, který je na (4,2) dle dohody. Jak jsi na tom ty? Můžeme ukončit? hotovo?} \\
            }
        }
    }
     \\ \\

    \theutterance \stepcounter{utterance}  
    & & & \multicolumn{2}{p{0.3\linewidth}}{
        \cellcolor[rgb]{0.9,0.9,0.9}{
            \makecell[{{p{\linewidth}}}]{
                \texttt{\tiny{[GM$|$GM]}}
                \texttt{ŘEKNI: Moje aktuální pozice: P (1,1), O (1,3), T (4,2), S (5,5), M (5,7), A (7,1), D (7,7). Všechny objekty mám na cílových pozicích kromě T, který je na (4,2) dle dohody. Jak jsi na tom ty? Můžeme ukončit? hotovo?} \\
            }
        }
    }
    & & \\ \\

    \theutterance \stepcounter{utterance}  
    & & \multicolumn{4}{p{0.6\linewidth}}{
        \cellcolor[rgb]{0.9,0.9,0.9}{
            \makecell[{{p{\linewidth}}}]{
                \texttt{\tiny{[P1$\langle$GM]}}
                \texttt{Úspěšně přesunuto 'D' do (7,7). Vaše aktualizovaná mřížka vypadá takto:} \\
\texttt{```} \\
\texttt{ 1234567} \\
\texttt{╔═╤═════╗} \\
\texttt{║P│◌◌◌◌A║ 1} \\
\texttt{╟─┤◌T◌◌◌║ 2} \\
\texttt{║O└─┐◌◌◌║ 3} \\
\texttt{║◌◌◌├─┐◌║ 4} \\
\texttt{║◌◌◌│S│◌║ 5} \\
\texttt{╟─┬─┴─┼─╢ 6} \\
\texttt{║◌│◌◌M│D║ 7} \\
\texttt{╚═╧═══╧═╝} \\
\texttt{```} \\
\texttt{Právě hrajete kolo 15 z maximálně 28.} \\
\texttt{Dohromady jste nasbírali 24 z 16 pokut.} \\
\texttt{Druhý hráč poslal tuto zprávu:} \\
\\ 
\\ 
\texttt{"Moje aktuální pozice: P (1,1), O (1,3), T (4,2), S (5,5), M (5,7), A (7,1), D (7,7). Všechny objekty mám na cílových pozicích kromě T, který je na (4,2) dle dohody. Jak jsi na tom ty? Můžeme ukončit? hotovo?"} \\
\\ 
\\ 
\texttt{Jaký je váš další příkaz?} \\
            }
        }
    }
    & \\ \\

    \theutterance \stepcounter{utterance}  
    & \multicolumn{4}{p{0.6\linewidth}}{
        \cellcolor[rgb]{0.9,0.9,0.9}{
            \makecell[{{p{\linewidth}}}]{
                \texttt{\tiny{[P1$\rangle$GM]}}
                \texttt{ŘEKNI: Moje aktuální pozice: P (1,1), O (1,3), T (4,2), S (5,5), M (5,7), A (7,1), D (7,7). Všechny objekty sedí na cílových pozicích. hotovo!} \\
            }
        }
    }
    & & \\ \\

    \theutterance \stepcounter{utterance}  
    & & & \multicolumn{2}{p{0.3\linewidth}}{
        \cellcolor[rgb]{0.9,0.9,0.9}{
            \makecell[{{p{\linewidth}}}]{
                \texttt{\tiny{[GM$|$GM]}}
                \texttt{ŘEKNI: Moje aktuální pozice: P (1,1), O (1,3), T (4,2), S (5,5), M (5,7), A (7,1), D (7,7). Všechny objekty sedí na cílových pozicích. hotovo!} \\
            }
        }
    }
    & & \\ \\

    \theutterance \stepcounter{utterance}  
    & & & \multicolumn{2}{p{0.3\linewidth}}{
        \cellcolor[rgb]{0.9,0.9,0.9}{
            \makecell[{{p{\linewidth}}}]{
                \texttt{\tiny{[GM$|$GM]}}
                \texttt{true} \\
            }
        }
    }
    & & \\ \\

    \theutterance \stepcounter{utterance}  
    & & & \multicolumn{2}{p{0.3\linewidth}}{
        \cellcolor[rgb]{0.9,0.9,0.9}{
            \makecell[{{p{\linewidth}}}]{
                \texttt{\tiny{[GM$|$GM]}}
                \texttt{Initial states:} \\
\texttt{```} \\
\texttt{Player 1:} \\
\texttt{ 1234567} \\
\texttt{╔═╤═════╗} \\
\texttt{║◌│◌A◌◌◌║ 1} \\
\texttt{╟─┤◌T◌◌◌║ 2} \\
\texttt{║O└─┐◌◌◌║ 3} \\
\texttt{║◌◌◌├─┐◌║ 4} \\
\texttt{║◌◌S│◌│◌║ 5} \\
\texttt{╟─┬─┴─┼─╢ 6} \\
\texttt{║P│◌DM│◌║ 7} \\
\texttt{╚═╧═══╧═╝} \\
\\ 
\\ 
\texttt{Player 2:} \\
\texttt{ 1234567} \\
\texttt{╔═╤═════╗} \\
\texttt{║P│S◌D◌◌║ 1} \\
\texttt{╟─┤◌◌◌◌◌║ 2} \\
\texttt{║◌└─┐◌A◌║ 3} \\
\texttt{║◌◌M├─┐◌║ 4} \\
\texttt{║◌T◌│◌│O║ 5} \\
\texttt{╟─┬─┴─┼─╢ 6} \\
\texttt{║◌│◌◌◌│◌║ 7} \\
\texttt{╚═╧═══╧═╝} \\
\\ 
\texttt{```} \\
            }
        }
    }
    & & \\ \\

    \theutterance \stepcounter{utterance}  
    & & & \multicolumn{2}{p{0.3\linewidth}}{
        \cellcolor[rgb]{0.9,0.9,0.9}{
            \makecell[{{p{\linewidth}}}]{
                \texttt{\tiny{[GM$|$GM]}}
                \texttt{End states:} \\
\texttt{```} \\
\texttt{Player 1:} \\
\texttt{ 1234567} \\
\texttt{╔═╤═════╗} \\
\texttt{║P│◌◌◌◌A║ 1} \\
\texttt{╟─┤◌T◌◌◌║ 2} \\
\texttt{║O└─┐◌◌◌║ 3} \\
\texttt{║◌◌◌├─┐◌║ 4} \\
\texttt{║◌◌◌│S│◌║ 5} \\
\texttt{╟─┬─┴─┼─╢ 6} \\
\texttt{║◌│◌◌M│D║ 7} \\
\texttt{╚═╧═══╧═╝} \\
\\ 
\\ 
\texttt{Player 2:} \\
\texttt{ 1234567} \\
\texttt{╔═╤═════╗} \\
\texttt{║P│◌◌◌◌A║ 1} \\
\texttt{╟─┤◌T◌◌◌║ 2} \\
\texttt{║O└─┐◌◌◌║ 3} \\
\texttt{║◌◌◌├─┐◌║ 4} \\
\texttt{║◌◌◌│S│◌║ 5} \\
\texttt{╟─┬─┴─┼─╢ 6} \\
\texttt{║◌│◌◌M│D║ 7} \\
\texttt{╚═╧═══╧═╝} \\
\\ 
\texttt{```} \\
            }
        }
    }
    & & \\ \\

    \theutterance \stepcounter{utterance}  
    & & & \multicolumn{2}{p{0.3\linewidth}}{
        \cellcolor[rgb]{0.9,0.9,0.9}{
            \makecell[{{p{\linewidth}}}]{
                \texttt{\tiny{[GM$|$GM]}}
                \texttt{* success: True} \\
\texttt{* lose: False} \\
\texttt{* aborted: False} \\
\texttt{{-}{-}{-}{-}{-}{-}{-}} \\
\texttt{* Shifts: 13.00} \\
\texttt{* Max Shifts: 12.00} \\
\texttt{* Min Shifts: 6.00} \\
\texttt{* End Distance Sum: 0.00} \\
\texttt{* Init Distance Sum: 32.45} \\
\texttt{* Expected Distance Sum: 29.33} \\
\texttt{* Penalties: 24.00} \\
\texttt{* Max Penalties: 16.00} \\
\texttt{* Rounds: 15.00} \\
\texttt{* Max Rounds: 28.00} \\
\texttt{* Object Count: 7.00} \\
            }
        }
    }
    & & \\ \\

\end{supertabular}
}

\end{document}
